\chapter{Source formulas and common algebra}
\label{ch:cert-1}
Formula tags are local to their certificate section. All source and testing fields remain separately labelled.

\section{\texorpdfstring{A uniform native C1 cap and an ordinary-angle admission}{A uniform native C1 cap and an ordinary-angle admission}}
\label{component:p42ee07a067c1}

\label{p42ee07a067c1:sec:c1-global-native-cap}

This section uses the original canonical C1 source, not an enlarged
formal parameter family. Thus $\lambda=n=0$, $d>0$, $m,C,A,\alpha>0$,
$D,\beta\ge0$ with at most one zero, and the source has its actual
finite noncommuting stationary contact. Use the inherited global sign
entry and normalize $C=2$. The exact full-parent Alice reflection
permits $0<k\le1$. Put
\[
 x=k^2,\quad r=1-a\eta,\quad\theta=t\zeta,\quad
 h=1-\theta,\quad y=v_c^2=x\eta\zeta,
\]
\[
 K=1+x,\ R=1+xr,\ T=1+xh,\ G=h+\zeta,
 \quad M=R-t,\ J=T-a,\ E=1-x\eta h,\ Q=1-x\zeta r.
 \tag{NC1}
\]
The inherited exact signs give $0<r,\theta<1$, $E,Q,M,J>0$ and
$0<a,t<2$. There is no size assumption on $\eta,\zeta,e$ or a
positive lower bound on any opening.

\paragraph{Three short positive source identities.}
The original coquotient polynomials obey
\begin{align*}
 P_4&=TE+x(1-r)G,\\
 P_\alpha&=Rh(1+y)+\zeta ME,\\
 P_d&=p_0+\eta p_1,\qquad
 p_0=1+2xh+rx^2G+x\zeta,\quad
 p_1=x\{K\zeta+xh\theta\}.
 \tag{NC2}
\end{align*}
These are obtained by substituting $a\eta=1-r$ and $t\zeta=\theta$
in the unchanged seven native coquotients. In particular $P_d\ge1$.
The original nonnegative $D,\beta$ are precisely the inequalities
\[
 B_5:=K(1+y)(r+\eta)-(\eta T-1+r)P_d\ge0,
 \qquad B_\beta:=Kx(r+\eta)G-\theta P_d\ge0.
 \tag{NC3}
\]
The exact relative coefficients are
\begin{align*}
 \frac C d&=\frac{K(r+\eta)E}{2(1-r)P_d},&
 \frac A d&=\frac{\eta P_4}{2rP_d},\\
 \frac m d&=\frac{Kx(r+\eta)h(1+y)}{2M\zeta P_d},&
 \frac\alpha d&=\frac{x^2\zeta(r+\eta)P_\alpha}{2QP_d},\\
 \frac Dd&=\frac{B_5}{2J\eta P_d},&
 \frac\beta d&=\frac{B_\beta}{2\theta P_d}.
 \tag{NC4}
\end{align*}
All denominators have their original positive signs; no denominator is
silently declared a unit.

\begin{lemma}[A uniform native gap]
On every physical native C1 contact oriented by $0<k\le1$,
\[
                         E\ge\frac1{18}.             \tag{NC5}
\]
This includes the literal boundaries $D=0$ or $\beta=0$.
\end{lemma}
\begin{proof}
Define
\[
 \mathcal A=\theta T+K\zeta>0,\qquad
 \mathcal C=RG(1+rxh),
\]
\begin{align*}
 \mathcal B={}&h^2rx^2+hrx(2-h)+h^2x+hrx^2\zeta
             +hx\zeta+hx+h\\
             &+rx\zeta+x\zeta+2\zeta+1.
\end{align*}
Direct multiplication of (NC2)--(NC3), using $x\eta h=1-E$,
gives the polynomial identity
\[
                  hB_5=-\mathcal A E^2+\mathcal B E-\mathcal C.
 \tag{NC6}
\]
Every summand in $\mathcal B$ is nonnegative. Since
$0<x,r,h\le1$ and $h(2-h)\le1$, its six terms without $\zeta$
sum to at most $6$ and the coefficients of $\zeta$ sum to at most
$6$. Hence $\mathcal B\le6(1+\zeta)$. The native inequality
$\theta=t\zeta<2\zeta$ gives $1<h+2\zeta$, so
$1+\zeta<3G$. Also $\mathcal C\ge G>0$. From $B_5\ge0$ and
(NC6),
\[
 G\le\mathcal C\le\mathcal B E\le18GE.
\]
Cancel the positive pure source $G$. This proves (NC5) without
inverting an opening or assuming any of $\eta,\zeta$ bounded.
\end{proof}

\paragraph{An elementary division-free elimination.}
Keep $p_0,p_1$ from (NC2), and set
\begin{align*}
 a_2&=Kx\zeta-Tp_1,&
 a_1&=K(1+x\zeta r)-Tp_0+(1-r)p_1,&
 a_0&=Kr+(1-r)p_0,\\
 b_1&=KxG-\theta p_1,&
 b_0&=KxrG-\theta p_0.
\end{align*}
Then $B_5=a_2\eta^2+a_1\eta+a_0$ and
$B_\beta=b_1\eta+b_0$. Expanding these two expressions, or just
eliminating their quadratic and linear terms successively, gives
\begin{align*}
 &b_1^2B_5-a_2B_\beta^2-(b_1a_1-2a_2b_0)B_\beta\\
 &\qquad=b_1^2a_0-b_1a_1b_0+a_2b_0^2\\
 &\qquad=x^2h^2K^2RG^2
       \{1+2xh+rx^2h^2+(1-r)x\zeta\}.
 \tag{NC7}
\end{align*}
The last equality is a short polynomial identity in $x,r,h,\zeta$;
all five of its input coefficients are printed above. No inverse of
$b_1$, no resultant premise, and no localization is used. For reference,
$b_1=xh\{K(1+\zeta)-x\theta^2\}>0$ as well.

\begin{theorem}[Actual unit/unit C1 contacts enter the five-event exit]
If both actual C1 half-angles $k,v_c$ are positive bounded units,
the original source has five consecutive events bounded below relative
to $d$. Consequently its entire original parent has the inherited
all-tests tensor-SOS certificate. All upper event scales and the
retained zero-coefficient boundaries are allowed.
\end{theorem}
\begin{proof}
After the exact Alice orientation, $x,y$ are bounded positive units.
Thus $\eta\zeta=y/x$ is a bounded positive unit. If $\eta$ were
unbounded, $\zeta$ would be infinitesimal. Then $\theta=t\zeta$
would be infinitesimal, contradicting $x\eta(1-\theta)<1$.
If $\zeta$ were unbounded, $\eta$ would be infinitesimal and
$r=1-a\eta$ would have residue one, contradicting $x\zeta r<1$.
Therefore $\eta,\zeta$ are both bounded positive units.

All polynomials in (NC2)--(NC3) now have bounded coefficients and
bounded arguments. If $h$ were infinitesimal, the exact specialization
\[
             B_\beta\big|_{\theta=1}
                    =-1-x\zeta(1-r)\le-1
\]
would contradict $B_\beta\ge0$. Thus $h$ is a positive bounded
unit. The uniform cap gives that $E$ is a unit. Equations (NC2) make
$P_4,P_\alpha,P_d$ positive bounded units: use $P_4\ge E$,
$P_\alpha\ge Rh(1+y)$, and $P_d\ge1$.

The first four equations of (NC4) now show that
$C,A,m,\alpha$ are all strong relative to $d$. Their denominators
are positive bounded quantities; they may be infinitesimal. Suppose
both $D/d$ and $\beta/d$ were infinitesimal, including a possible
literal zero. Their last two bounded denominator products in (NC4)
would make $B_5,B_\beta$ infinitesimal. Every coefficient on the
left of (NC7) is bounded, so that side would be infinitesimal. Its
right side is a positive unit, since $x,h,K,R,G$ are units and the
last bracket is at least one. This is a contradiction.

Thus $D$ or $\beta$ is strong. The original cycle
$(\lambda,D,C,A,\alpha,m,\beta,n)$ has respectively the five
consecutive strong events $(D,C,A,\alpha,m)$ or
$(C,A,\alpha,m,\beta)$. Apply the inherited five-consecutive-event
theorem, whose stated scope allows arbitrary upper weights and every
independent testing chart. No independent test is identified with
an actual source angle in this source-entry argument.
\end{proof}



\section{\texorpdfstring{An exact beta-shared normalization for the two eta-large C1.4 charts}{An exact beta-shared normalization for the two eta-large C1.4 charts}}
\label{component:p10ee3d96a3f1}

\label{p10ee3d96a3f1:sec:r8-c1-large-eta-source-normalization}
This section proves an exact source reduction, not an all-tests admission.
It applies simultaneously to C14.1 and C14.2: the original native C1
source has $x=k^2$ positive infinitesimal, $L=x\zeta$ bounded below by a
positive real, and $\eta$ unbounded. All original physicality and
residual-complement conditions are retained. We normalize $d=1$.
The conclusion includes every positive source opening and the literal
faces $D=0$ and $\beta=0$. No testing variable occurs in the source
normalization.

\paragraph{A uniform strengthening of the source gap.}
Write $\sigma=\eta^{-1}$, $\theta=1-h$, $K=1+x$, and
$e=xh/\sigma=1-E$. The exact coefficients in (LH1), with no size
specialization, give
\[
 B_5=-xh\{KL+x\theta(1+xh)\}\eta^2+x b_1\eta+a_0,
\]
where
\begin{align*}
 b_1={}&-Lhrx-Lh-Lr+L-h^2rx^2+h^2rx-3h^2x\\
 &-2hrx+hx-3h+1,\\
 a_0={}&Kr+(1-r)\{1+L+2xh+rxL+rx^2h\}.
\end{align*}
Since $0<r,h<1$, $Lr<1$, and $L$ is bounded below,
$b_1/L$ and $a_0/L$ are bounded uniformly in all openings; the second
is positive. Dividing the inequality $B_5\ge0$ by $L\eta>0$ gives
\[
 e\{K+x\theta(1+xh)/L\}
 \le x(b_1/L)+\sigma(a_0/L)\le O(x)+O(\sigma).
 \tag{BE1}
\]
The bracket is at least one. Therefore $e\in\Inf_{>0}$ and $E\sim1$
on both eta-large charts. This proof does not assume $x/\sigma$ bounded,
and does not infer that $xh/\sigma^2$ is bounded.

\paragraph{Two exact identities and their bounded coefficients.}
Set
\[
 z=L^{-1},\quad b=Lr,\quad Q=1-b,\quad
 R=1+xr,\quad M=R-xz\theta,\quad J=1+xh-\sigma(1-r).
\]
Here $z>0$ is bounded, $0<b,r<1$, and $r=bz$.
The opening-cleared denominator is
\[
 P=\sigma zP_d
 =\sigma+K+\sigma z(1+bx+2xh)
   +\sigma bx^2hz^2+x^2h\theta z.
\]
In particular $P,M,J\sim1$. Introduce
\[
 H_0=1+z-r=1+zQ,\qquad B=P+xz(J-\sigma)
\]
and
\[
 \mathcal H=RH_0+z\left[(1-r)xh(2+xr)
 -e\{2+x(1+h+r\theta)\}+Ke^2\right].
 \tag{BE2}
\]
Then $B\sim1$, and $\mathcal H-H_0\in\Inf$, while $H_0\ge1$ is
bounded. Thus $B$ and $\mathcal H$ are positive bounded units.
The original NC coefficients satisfy exactly
\begin{align*}
 2\theta P\beta&=hB-\sigma H_0,                    \tag{BE3}\\
 2P(JD+x\theta\beta)&=\sigma^2\mathcal H.         \tag{BE4}
\end{align*}
To verify these without any cancellation of an opening, use
\[
 N_D=K(1+\sigma z)(1+\sigma r)-JP=2JPD,
 \quad
 N_\beta=K(1+\sigma r)(1+xhz)-\theta P=2\theta P\beta.
\]
Expansion gives $N_\beta=hB-\sigma H_0$. The identity
$N_D+xN_\beta=\sigma^2\mathcal H$ follows on inserting
$\sigma e=xh$. These are polynomial identities after multiplication
by $\sigma^2$; no resultant, opening inverse, or approximation is used.
As an immediate consequence,
\[
 D=O(\sigma^2),\qquad x\theta\beta=O(\sigma^2),\qquad
 JD+x\theta\beta\asymp\sigma^2.
 \tag{BE5}
\]
In particular the collapsing $D$ pivot cannot be given a fixed positive
unit cap in this source chart. The estimate is stronger than bounding
$D$ by $O(\sigma^2+x\sigma)$ before using original beta nonnegativity.

\paragraph{A source sphere retaining both coefficient boundaries.}
All radicals below are pure source. Define
\[
 d_0=\sqrt{\frac{2PJD}{\sigma^2\mathcal H}},\qquad
 b_0=\sqrt{\frac{2Px\theta\beta}{\sigma^2\mathcal H}}.
\]
Then exactly
\[
 d_0,b_0\ge0,\quad d_0^2+b_0^2=1,
 \quad D=\frac{\sigma^2\mathcal H}{2PJ}d_0^2,
 \quad \beta=\frac{\sigma^2\mathcal H}{2Px\theta}b_0^2.
 \tag{BE6}
\]
No inverse of $D$, $\beta$, $d_0$, or $b_0$ occurs. In particular
$D=0$ corresponds to $(d_0,b_0)=(0,1)$, and $\beta=0$ to $(1,0)$.
Equation (BE3) now becomes the exact positive decomposition
\[
 \frac h\sigma=\frac{H_0}{B}
       +\frac\sigma x\frac{\mathcal H}{B}b_0^2.
 \tag{BE7}
\]
It displays, rather than discards, the possibly unbounded factor
$h/\sigma$. The latter is bounded below by a positive unit, but may
be unbounded when $\sigma/x$ is unbounded.

Put
\[
 \kappa=\frac{K(1+\sigma r)(1+\sigma z)}{2MP}\sim\frac12,
 \qquad m_0=\kappa x^2z\frac{H_0}{B},\qquad
 m_1=\kappa x\sigma z\frac{\mathcal H}{B}.
\]
The original $m$ coefficient has the exact split
\[
 m=m_0+m_1b_0^2,\qquad
 m_0\asymp x^2z,\quad m_1\asymp x\sigma z.
 \tag{BE8}
\]
Thus the same nonnegative source coordinate $b_0$ occurs in the beta
packet and in the extra part of $m$. This is an exact beta-aware
replacement for an assumed bounded $\mu=h/\sigma$.
All factors hidden in the two comparisons in (BE8) are bounded positive
units uniformly in the original openings.
For the other three positive coefficients the opening-cleared products
are also units:
\[
 \frac{C(1-r)}z\sim\frac12,\qquad Ar\sim\frac{H_0}{2},
 \qquad \alpha Qz\sim\frac12.
 \tag{BE9}
\]
Indeed their exact expressions are those in the original large-$L$
source chart, which are algebraic for bounded $z$ as well; use (BE1).
For example $Ar=[z(1+xh)E+(1-r)(1+xhz)]/(2P)$.
No lower bound on $r$, $1-r$, $Q$, or $\theta$ was imposed.

\paragraph{Exact full-parent preservation.}
For arbitrary independent reciprocal testing coordinates $u,w$, put
$D_u=1+u^2$, $D_w=1+w^2$ and
\begin{gather*}
 M_w=W-wX+uY-uwZ,\quad R_w=uW-Y,\\
 J_w=uwW+uX-wY-Z,\quad B_w=X+uZ,\\
 K_w=J_w^2+D_u(Y-wZ)^2+D_wB_w^2-D_uD_wZ^2.
\end{gather*}
The entire original reciprocal parent is exactly
\begin{align*}
 F_w={}&4m_0M_w^2+4m_1(b_0M_w)^2
       +4CD_wR_w^2+4AD_uD_wY^2\\
 &+\frac{2\sigma^2\mathcal H}{PJ}(d_0J_w)^2
       +4\alpha D_u(Y-wZ)^2
       +\frac{2\sigma^2\mathcal H}{Px\theta}D_w(b_0B_w)^2
       +2K_w.                                      \tag{BE10}
\end{align*}
This is an equality in the complete four-state quadratic form.
It therefore preserves every first state normal and every hard
quadratic coupling; it is not a restriction to a contact line or a
leading form. Every new coefficient or radical is pure source, and all
$D_u,D_w$ factors expand as finite testing sums of squares.
The identity is polynomial in $u,w,W,X,Y,Z$ and therefore includes all
signs and every independent testing rate. It holds at literal $u=0$
and $w=0$. Bihomogenizing $u$ and $w$ gives the corresponding identity
at their other literal projective endpoints, without division by a
testing coordinate. Applying the real identity to real and imaginary
state parts gives the complex-state identity as well.

The source sphere, (BE7), and (BE8) do not make $2K_w$ nonnegative.
They leave a precise full-parent allocation problem: use the common
$b_0$ packet in (BE10) together with the $d_0$ packet, uniformly through
$d_0=0$ and $b_0=0$. Consequently neither C14.1 nor C14.2 is admitted
by this section. The sigma-unit six-root proof's positive-unit $D$
pivot cannot be imported unchanged. No claim of nonemptiness of any
residual-complement source domain is made.

\paragraph{Verification scope.}
The accompanying narrow checker verifies (BE1)'s polynomial input,
(BE3), (BE4) in both reciprocal and native notation, the exact $m$
split, the opening-cleared coefficient products, and (BE10) as the
complete quadratic polynomial, including its bihomogenization.
The asymptotic deductions use the ordered-field inequalities above;
checker success alone is not an all-tests certificate.

\paragraph{Targeted full-normal extension and the exact null direction.}
Here is a test of the proposed normalization on the normal block, with
no source opening omitted. At $w=0$, put $X=N_x-uhZ$ and $t=u^2$.
The original parent splits into its manifestly nonnegative $(W,Y)$
block and the following complete $(N_x,Z)$ block:
\begin{align*}
 Q_{xz}={}&[4Dt+4\beta+2(1+t)]N_x^2\\
 &+4u[-2D(1+th)+2\beta\theta-h(1+t)]N_xZ\\
 &+[4D(1+th)^2+4\beta t\theta^2+2t(1+t)h^2]Z^2.
 \tag{BE11}
\end{align*}
At the literal endpoint $u=w=0$ this is $(4\beta+2)X^2+4DZ^2$.
Thus when $D=0$, the vector $(W,X,Y,Z)=(0,0,0,1)$ is an
\emph{exact} null direction of the full parent. The old six-root
argument cannot impose a strictly positive pivot on $\sigma N_z$
at this face. This is an obstruction to that proof architecture,
not to positivity; $w=0$ itself already has an elementary SOS.

The sphere suggests the legal replacement roots
\[
 \Gamma=\sigma Z,\qquad Z_d=d_0\Gamma,\qquad
 Z_b=b_0(u/k)\Gamma,\qquad Z_e=\sqrt e\,(u/k)\Gamma.
\]
Set $p=\sigma b_0/k$. Equation (BE7) implies $p$ bounded, since
$hB=\sigma H_0+\mathcal H p^2$. Also
\[
 \frac{xh^2}{\sigma^2}
   =\frac{eH_0}{B}+\frac{\mathcal Hh}{B}b_0^2.
\]
Consequently the entire $Z^2$ line of (BE11) equals
\[
 \frac{2\mathcal H}{PJ}(1+th)^2Z_d^2
 +2\mathcal H\left\{\frac\theta P+\frac{(1+t)h}{B}\right\}Z_b^2
 +\frac{2(1+t)H_0}{B}Z_e^2.                       \tag{BE12}
\]
For $u$ infinitesimal each displayed coefficient has a positive-unit
residue, including $\theta\to0$, because $\theta+h=1$.
To retain the first normal exactly, define
\[
 A_x=\frac{H_0}{B}\left\{\frac{xz(J-\sigma)}P-t\right\}
       -\frac{H_0}P,
\]
\begin{align*}
 c_d&=4ud_0\left\{A_x-\frac{\sigma\mathcal H}{PJ}(1+th)\right\},\\
 c_b&=4kb_0\left\{A_x+\frac{\mathcal H}{B}
                       \frac{\sigma z(J-\sigma)}P\right\}
                      -4tp\frac{\mathcal H}{B}.
\end{align*}
The middle line of (BE11) is exactly
$c_dN_xZ_d+c_bN_xZ_b$. Both coefficients are infinitesimal when
$u$ is infinitesimal, through every source opening and either sphere
endpoint. Extracting $4(\beta-\beta_*)N_x^2$ for any bounded
$0\le\beta_*\le\beta$ leaves a bounded Gram on
$(N_x,Z_d,Z_b,Z_e)$ with strictly positive diagonal residue. Hence
this normal block has the standard bounded positive-Gram lift.
No inverse of $d_0,b_0,u$, or $\theta$ enters these coefficients.
At $d_0=0,u=0$ all three $Z$ roots vanish, as the exact null direction
requires. At $b_0=0$, the $Z_b$ root and its cross coefficient vanish.

For the full $w\ne0$ contact conductor the corresponding proposed
roots are $d_0(\sigma Z-uwW)$,
$b_0(u/k)(\sigma Z-uwW)$, and
$\sqrt e\,(u/k)(\sigma Z-uwW)$, together with the existing radial
root and angular roots. Equations (BE11)--(BE12) verify the formerly
collapsed normal direction and its first-normal coupling, but do not
allocate the scalar contact polynomial and the $w$-dependent radial
couplings. Their \emph{joint} allocation is the remaining step for
this proposed frame. No positivity claim for that uncomputed Gram is
made here, and no endpoint already known to be elementary is counted
as a new source-domain admission.


\section{\texorpdfstring{Shared event mechanisms for the remaining C1 sources}{Shared event mechanisms for the remaining C1 sources}}
\label{component:p89fd3edb3edd}

\label{p89fd3edb3edd:sec:r6-c1-common-remote}
Normalize the original parent by its positive pure-source coefficient $d$.
Write $D_z=1+z^2$, and write the signed core as $2K$, where
\[
 K=(uW+uvX-Y-vZ)^2+D_u(vY-Z)^2+D_v(X+uZ)^2-D_uD_vZ^2.
 \tag{CE1}
\]
The event packets used below are
\[
4m(vW-X+uvY-uZ)^2,\quad4CD_v(uW-Y)^2,
\quad4AD_uD_vY^2,\quad4\alpha D_u(vY-Z)^2.
\]
Every omitted event and every nonnegative excess is retained as its exact
original square. Source and testing fields remain independent. ``Strong''
means bounded below by a positive real number; ``unbounded'' means larger
than every real number. Neither assertion constrains upper scales.

\begin{theorem}[Two adjacent unbounded events, with complete raw frames]
If $C,A$ are unbounded, the full original parent is tensor SOS at every
projective test for which $u$ is not infinitesimal, and every state.
No physicality assumption beyond nonnegative event weights is needed for
this particular assertion. By the exact cyclic relabeling below, if
$A,\alpha$ are unbounded, the conclusion holds whenever $v$ is not
projectively unbounded.
\end{theorem}
\begin{proof}
Select $C=A=\mu$, with a positive real $\mu$ chosen after the real testing
residues; unboundedness permits every such choice. The following are exact
full quadratic forms, before any restriction to a state ray or contact.
They therefore retain every first state normal and hard-normal coupling.

For $u$ a bounded nonzero unit and $v$ bounded, use the four raw roots
\[
 (W,Y,B,V),\qquad B=X+uZ,\quad V=vZ.
\]
The entire core equals
\[
 K=(uW-Y+uvB)^2+D_vB^2+D_uv^2Y^2+D_uu^2V^2
       -2D_uuV(W+vB).                         \tag{CE2}
\]
The two selected event forms remain $4\mu D_v(uW-Y)^2$ and
$4\mu D_uD_vY^2$. In their common real kernel $W=Y=0$, the core has
matrix on $(B,V)$
\[
 \begin{pmatrix}1+v^2D_u&-D_uuv\\-D_uuv&D_uu^2\end{pmatrix},
 \qquad \det=D_uu^2>0.                       \tag{CE3}
\]
Here and in residue matrices, bounded testing letters mean their real
residues. Thus a sufficiently large finite $\mu$ gives a strictly positive
residue on the entire four-root space, even when $v=0$ literally.

For $p=1/u$ infinitesimal and $v$ bounded, homogenize by $p^2$ and take
\[
 (W,Y,B,R),\qquad B=pX+Z,\quad R=vX.
\]
The exact full core is
\[
 K_p=B^2-2pvBW-2vBY+D_pR^2+2D_pRW
                  +(W-pY)^2+D_pv^2Y^2.       \tag{CE4}
\]
The selected packets are $4\mu D_v(W-pY)^2$ and
$4\mu D_pD_vY^2$. At $p=0$ their sum with $2K_p$ is
\[
 4\mu D_v(W^2+Y^2)+2(W+R)^2+2(B-vY)^2,       \tag{CE5}
\]
strict for every $\mu>0$. The formula allows literal $p=0$ and $v=0$;
the raw roots need not form independent state coordinates.

For projectively unbounded $v$, put $w=1/v$, homogenize by $w^2$, and
use $(W,X,Y,H)$ with $H=wZ$. The exact full core is
\begin{align*}
 K_w={}&u^2H^2-2u^2HY-2uHW+2uwHX+u^2w^2W^2\\
 &+2u^2wWX-2uw^2WY+(D_u+w^2)(X^2+Y^2)-2uwXY.
                                                        \tag{CE6}
\end{align*}
The selected packets are $4\mu D_w(uW-Y)^2$ and
$4\mu D_uD_wY^2$. If $u$ is a bounded nonzero unit, at $w=0$ their
common kernel has $W=Y=0$ and $K_w=D_uX^2+u^2H^2$ strictly positive.
A sufficiently large real $\mu$ again works. If also $p=1/u$ is
infinitesimal, homogenize (CE6) by $p^2$. The exact full core is
\begin{align*}
 K_{pw}={}&H^2-2pHW+2pwHX-2HY+w^2W^2+2wWX-2pw^2WY\\
 &+(D_p+p^2w^2)(X^2+Y^2)-2pwXY.              \tag{CE7}
\end{align*}
At $p=w=0$ the full selected residue is
\[
 4\mu(W^2+Y^2)+2X^2+2(H-Y)^2>0.             \tag{CE8}
\]
All these exact Gram coefficients are bounded. Strictness in (CE3) and
its counterpart following (CE6) implies a finite cap by the elementary
real Schur-complement argument: the event Gram is positive definite on
the complement of its kernel, and the core is positive definite on the
kernel. No source/testing rate comparison is involved. Apply the bounded
positive-residue Gram lemma, and restore every retained original square.

For the second assertion, compose the inherited Alice reflection with
the swap/reflection. The resulting one-step inverse cycle rotation is
\[
 (u',v')=(1/v,u),\quad
 (W',X',Y',Z')=(Y+vZ,-W-vX,Z-vY,vW-X),          \tag{CE9}
\]
and sends event index $i$ to $i-1$ in
$(\lambda,D,C,A,\alpha,m,\beta,n)$. In particular original $A,\alpha$
become new $C,A$. Its full-parent identity is
\[
 v^2F(\mathrm{rot}^{-1}c;1/v,u;\psi')=D_vF(c;u,v;\psi).
                                                        \tag{CE10}
\]
The state map has determinant of absolute value $D_v^2$, and its
homogeneous version at $v=\infty$ is invertible as well. Thus
$u'=1/v$ not infinitesimal is exactly $v$ not projectively unbounded.
Only pure testing factors are restored. Both signs, all literal endpoints,
and complex states (by real/imaginary splitting) are included.
\end{proof}

\begin{theorem}[Four consecutive strong events]
For an original physical parent with $C,A,\alpha,m$ strong, the full
parent is tensor SOS at every test except possibly $u$ a bounded nonzero
unit and $v$ projectively unbounded. All original upper coefficient
scales and all nonnegative unselected coefficients are allowed.
\end{theorem}
\begin{proof}
Cap these four events at any positive real $\mu$ smaller than all four
weights, and retain every excess. Noncommuting-residue tests use the
inherited physical-source theorem, applied to the original full parent.
For commuting-residue tests use the capped parent.

First suppose $v$ is bounded. On the three open commuting axes, the core
and the four caps have zero common kernel: at $u=0$, the core forces
$X=Y=0$, then $m$ and $\alpha$ force $W=Z=0$ when $v$ is a unit;
at $u=\infty$, $C,A$ force $W=Y=0$, then the core forces $X=Z=0$;
at $v=0$, $A,\alpha$ force $Y=Z=0$, then the core forces $W=X=0$
when $u$ is a unit. These are strictly positive bounded state Grams.

At $(u,v)=(0,0)$ the residue is
\[
 (2+4\mu)X^2+(2+8\mu)Y^2+4\mu Z^2.
\]
Only the $W$ ray remains. Normalize $W=1$ and take the five infinitesimal
roots $(u,v,X,Y,Z)$. The full polynomial has no constant or linear term;
its quadratic part is
\[
 2(u-Y)^2+2X^2+
 4\mu\{(v-X)^2+(u-Y)^2+Y^2+Z^2\},            \tag{CE11}
\]
which is strictly positive. At $(u,v)=(\infty,0)$ only the $X$ ray
remains. Normalize $X=1$ and use $(p,v,W,Y,Z)$. Its quadratic part is
\[
 2(W+v)^2+2(p+Z)^2+
 4\mu\{(p+Z)^2+W^2+Y^2+Z^2\},              \tag{CE12}
\]
also strictly positive. All higher monomials in each chart allocate to
products of two displayed roots with infinitesimal bounded coefficients.
This is an exact finite Gram, not a truncation certificate.

For $w=1/v$ infinitesimal use the full raw frame of (CE6). All four
selected packets, with every normal present, are exactly
\[
 4\mu\{(W-wX+uY-uH)^2+D_w(uW-Y)^2
              +D_uD_wY^2+D_u(Y-H)^2\}.        \tag{CE13}
\]
At $u=w=0$ their sum with $2K_w$ has residue
\[
 4\mu W^2+2X^2+(8\mu+2)Y^2+4\mu(Y-H)^2,    \tag{CE14}
\]
strict on all four roots. If $p=1/u$ is infinitesimal, homogenize again;
the exact packets become
\[
 4\mu\{(pW-pwX+Y-H)^2+D_w(W-pY)^2
              +D_pD_wY^2+D_p(Y-H)^2\}.        \tag{CE15}
\]
Together with (CE7), at $p=w=0$ the residue is
\[
 4\mu W^2+2X^2+4\mu Y^2+(8\mu+2)(Y-H)^2,    \tag{CE16}
\]
again strict. These frames require no state normalization, and no
division by either small testing coordinate. Literal $w=0$ is included.
The only chart not covered is ordinary $u$ and infinitesimal $w$.
Restore the original packets and all pure factors as above.
\end{proof}

\begin{lemma}[Matched weak outer events at the large-Bob corner]
Suppose $q>0$ is a pure-source infinitesimal and
\[
 C\ge c q^2,\quad m\ge m_*q^2,\quad
 A\ge a_*/q^2,\quad\alpha\ge b_*/q^2
                                                        \tag{CE17}
\]
for positive real $c,m_*,a_*,b_*$. At projectively unbounded $v$,
the complete parent is tensor SOS whenever $u$ is infinitesimal or
projectively unbounded. In the latter chart the lower bound on $m$
is unnecessary.
\end{lemma}
\begin{proof}
Keep $w=1/v$ and $H=wZ$ as in (CE6)--(CE7). If $u$ is infinitesimal,
use the five raw roots
\[
 E=qW,\quad X,\quad F=Y/q,\quad G=H/q,\quad V=uwW.
\]
The exact whole core (CE6) becomes
\begin{align*}
 K_w={}&V^2+(D_u+w^2)X^2
       +q^2\{u^2(G-F)^2+D_wF^2\}\\
 &-2uGE+2quwX(G-F)+2uVX-2uw^2EF.             \tag{CE18}
\end{align*}
Cap only $m,A,\alpha$ at the indicated weights. Their complete packets
are respectively
\[
 4m_*(E-qwX+uq^2(F-G))^2,\quad
 4a_*D_uD_wF^2,\quad4b_*D_u(F-G)^2.
\]
The selected residue is
\[
 2V^2+2X^2+4m_*E^2+4a_*F^2+4b_*(F-G)^2>0.   \tag{CE19}
\]
Every coefficient in (CE18) is bounded independently of all relative
source/testing rates.

If instead $p=1/u$ is infinitesimal, use
\[
 E=qW,\quad F=Y/q,\quad G=H/q,\quad
 R=X+wW,\quad S=pX.
\]
The exact whole core (CE7) is
\[
 K_{pw}=R^2+D_wS^2+q^2(G-F)^2+p^2D_wq^2F^2
       -2pE(G+w^2F)+2qwS(G-F).               \tag{CE20}
\]
Cap only $C,A,\alpha$. Their complete packets are
\[
 4cD_w(E-pq^2F)^2,\quad4a_*D_pD_wF^2,
 \quad4b_*D_p(F-G)^2.
\]
Together with $2K_{pw}$ their residue is
\[
 2R^2+2S^2+4cE^2+4a_*F^2+4b_*(F-G)^2>0.     \tag{CE21}
\]
Only the pure-source inverse $1/q$ occurs in a root. These exact raw
frames include every state normal and every literal endpoint. The
strict-residue lemma and restoration of excess packets prove the claim.
\end{proof}

\paragraph{Exact application to the original C1 source rows.}
The source premises here are precisely the already proved original
formulas (NC1)--(NC7), (NI3)--(NI10); no independent source family is
substituted. These common theorems give the following new quantified
test reductions, always intersected with prior admissions:
\begin{enumerate}
\item C1.1, every $e$ type: (NI3) gives $C,A$ unbounded;
only $u$ infinitesimal remains.
\item C1.4, $L$ infinitesimal and $\eta$ unbounded: (NI6)
again gives $C,A$ unbounded; only $u$ infinitesimal remains.
\item C1.4, $L$ unbounded: (NI4) and (CE20) leave only
$u$ bounded and $v$ unbounded.
\item C1.5 actual UH, $\eta$ unbounded and $\zeta$ unit:
(NI9) gives four strong events; only $u$ unit and $v$ unbounded remains.
\item C1.5 actual UH, $\zeta$ unbounded: (NI10) and (CE17)
leave the same sole strip, $u$ unit and $v$ unbounded.
\end{enumerate}
For the last line choose $q=\zeta^{-1/2}$. Formula (NI10) gives
$C,m\asymp q^2$, $\alpha\asymp1/(q^2Q)$ and $A\asymp1/r$;
the original $Q=1-x\zeta r>0$, with unit $x$, bounds $Aq^2$ below.
Thus (CE17) holds with positive real constants. The adjacent-event
theorem covers bounded $v$ and (CE18)--(CE21) remove both end types of
$u$ at unbounded $v$. Both actual UH source patterns now have the same
sole untreated testing strip, ordinary $u$ and unbounded $v$.

In the $L$-unbounded C1.4 line take $q=L^{-1/2}$.
Then (NI4) gives $C\asymp q^2$, $\alpha\asymp1/(q^2Q)$ and
$Aq^2\asymp1/(Lr)>1$. The $p$-chart (CE20), which needs no lower
bound on $m$, removes the projectively unbounded-$u$ tests there.

All bounds are relative to original $d$. The four-source-event result
does not require the additional NI9 total-scale condition that at least
one opening be infinitesimal. Every literal permitted $D=0$ or
$\beta=0$ is retained throughout, since neither coefficient was capped.
No claim of nonemptiness or full-row completion is made. In particular,
the $L$ unit and the $L$ infinitesimal/$\eta$ bounded parts of C1.4 are
unchanged, and these remote results do not supply the missing critical
contact conductors.



\section{\texorpdfstring{Further exact source information on the three broad C1 rows}{Further exact source information on the three broad C1 rows}}
\label{component:p51451360ff4a}

\label{p51451360ff4a:sec:c1-residual-source-implications}

This section records source implications of (NC1)--(NC7), not further
all-tests certificates. All quantities are in the coefficient field.
Write $f\asymp g$ for two strictly positive quantities whose quotient
and reciprocal quotient are bounded. Set
\[
             L=x\zeta,\qquad S=1+L+y,\qquad \rho=r+\eta.
\]
The native signs and (NC5) imply the uniform estimates
\[
 S\le P_d\le6S,\qquad
 \frac{1+L}{54}\le P_4\le3(1+L),\qquad
 P_\alpha\asymp h(1+y)+\zeta M.                 \tag{NI1}
\]
Indeed $rx^2\zeta<x$ follows from $Q>0$, and
$\eta x^2h\theta=x\theta(1-E)\le1$, giving the $P_d$ bounds
from (NC2). For $P_4$, if $L\le2$ use $E\ge1/18$; if $L\ge2$,
then $r<1/L\le1/2$ and $P_4\ge L/2$. Its upper bound follows
immediately from (NC2). The last comparison uses $1\le R\le2$
and $1/18\le E<1$.
Consequently the unchanged original source satisfies
\begin{align*}
 C/d&\asymp\frac{\rho}{(1-r)S},&
 A/d&\asymp\frac{\eta(1+L)}{rS},\\
 m/d&\asymp\frac{x\rho h(1+y)}{M\zeta S},&
 \alpha/d&\asymp
 \frac{x^2\zeta\rho\{h(1+y)+\zeta M\}}{QS}.
 \tag{NI2}
\end{align*}
No opening has been bounded away from zero in these statements.
In particular, $E$ and $P_4$ do not supply additional hidden small
source denominators on C1.1, C1.4 or C1.5.

\paragraph{C1.1: two unavoidable large weights and its actual angle.}
Here $x,\zeta$ are infinitesimal and $\eta$ is unbounded.
The inherited original-$D$ cap gives $q=x\eta$ bounded with
residue at most $1/2$. Hence $y=q\zeta$ is infinitesimal:
the actual angle type is necessarily $(I,I)$, independently of $e$.
The exact formulas give
\[
 C/d\asymp\frac{\eta}{1-r},\qquad
 A/d\asymp\frac{\eta}{r},\qquad
 m/d\asymp\frac{e^2}{M},\qquad
 \alpha/d\asymp xq\zeta,\qquad
 \beta/d+\frac12\asymp\frac{e^2}{t}.             \tag{NI3}
\]
Thus $C,A$ are unbounded relative to $d$, and $\alpha$ is weak.
The inherited bounded polynomial formula for $B_5$ gives $D/d=O(1/\eta)$,
so $D$ is weak too, including its literal zero. In (NI3), $\beta\ge0$
forces $t=O(e^2)$. Therefore when $e$ is infinitesimal, $t$ is
infinitesimal, $M$ is a unit and $m$ is weak. When $e$ is a unit,
$m$ is strong; when $e$ is unbounded, both $m,\beta$ are unbounded.
These are restrictions on the original C1.1 row, not new admissions.
Its positive infinitesimal $\delta_{\rm tot}$ condition is automatic
because $C/d,A/d$ are unbounded.

\paragraph{C1.4: the natural scale is $L=x\zeta$.}
Now $x$ is infinitesimal and $\zeta$ is noninfinitesimal.
The following implications follow uniformly through all retained
coefficient boundaries.
\begin{enumerate}
\item If $L$ is unbounded, $Q>0$ forces $r$ infinitesimal, and
$1-r=a\eta<2\eta$ makes $\eta$ noninfinitesimal. Hence $y=L\eta$
is unbounded: the actual type is $(I,H)$. Here $\zeta$ is unbounded,
$t=\theta/\zeta$ is infinitesimal, and $M$ is a unit. Since
$h y/\zeta=x\eta h<1$, (NC2) gives $P_\alpha\asymp\zeta$.
Equations (NI1)--(NI2) then yield
\[
 C/d\asymp L^{-1},\qquad
 A/d\asymp r^{-1},\qquad
 \alpha/d\asymp L/Q,\qquad
 m/d\asymp (x\eta h)/\zeta\le O(\zeta^{-1}).
 \tag{NI4}
\]
In particular $C,m$ are weak, $A,\alpha$ unbounded, and
$(m/d)/(C/d)=O(x)$ is infinitesimal. The original total-scale
intersection is automatic.

\item If $L$ is a bounded unit, $\zeta$ is unbounded, $M$ is a
unit, $P_\alpha\asymp\zeta$, and $S\asymp1+\eta$. The identity
$1=r+a\eta$ gives $\rho/(1+\eta)$ a positive bounded unit.
Thus
\[
 C/d\asymp(1-r)^{-1},\qquad
 \alpha/d\asymp Q^{-1},\qquad
 A/d\asymp\frac{\eta}{r(1+\eta)},\qquad
 m/d=O(x).
 \tag{NI5}
\]
So $C,\alpha$ are strong and $m$ weak. Moreover $A$ is weak
exactly when $\eta$ is infinitesimal; that case has $r$ a unit and
$D/d$ unbounded, by $B_5\sim1$ in (NC3). The actual Bob angle is
infinitesimal, a unit, or unbounded exactly as $\eta$ is, because
$y=L\eta$. No restriction on $D,\beta$ is added in the other cases.

\item If $L$ is infinitesimal and $\eta$ is bounded, then $y$ is
infinitesimal. The $B_\beta\ge0$ equation bounds
$\theta\le Kx\rho G/P_d=O(x+L)$, so $h,M$ are units.
Now (NI2) gives $C$ strong, $m,\alpha$ weak, and $A$ weak exactly
when $\eta$ is infinitesimal. In that last case $D/d$ is unbounded,
as in the preceding item. The actual type is $(I,I)$.

\item If $L$ is infinitesimal and $\eta$ is unbounded, then
\[
 C/d\asymp\frac{\eta}{(1-r)(1+y)},\qquad
 A/d\asymp\frac{\eta}{r(1+y)}.                    \tag{NI6}
\]
Both are unbounded since $\eta/(1+y)=(\eta^{-1}+L)^{-1}$.
Also $Q$ is a unit and (NI2) bounds
\[
 \alpha/d\le O\!\left(x^2\zeta\eta h+
              \frac{x^2\zeta^2\eta}{1+y}\right)=O(L),
\]
so $\alpha$ is weak. Finally $J\sim1$ and
$P_d/(1+y)\sim1$ directly in (NC2); (NC4) gives
$D/d+1/2\sim1/2$, hence $D$ weak. This extends the
$C,A$-large/$D,\alpha$-weak source pattern of C1.1 into C1.4,
without claiming that its existing conductors apply. Here $y$ may
be infinitesimal, a unit, or unbounded. The total-scale intersection
is again automatic.
\end{enumerate}
These four alternatives exhaust C1.4, but they are necessary source
reductions only. In particular no entire C1.4 leaf is closed.

\paragraph{C1.5 after the actual unit/unit admission.}
Here $k$ is a unit. The new theorem removes actual $(U,U)$; the
inherited both-$\eta,\zeta$-infinitesimal theorem is also subtracted.
Two useful remaining source types can be specified more tightly.

If $v_c$ is infinitesimal, $\eta\zeta$ is infinitesimal. Neither
$\eta$ nor $\zeta$ can be unbounded: $\zeta$ infinitesimal forces
$\theta$ infinitesimal and excludes unbounded $\eta$ by $E>0$;
$\eta$ infinitesimal gives $r\sim1$ and excludes unbounded $\zeta$
by $Q>0$. Thus exactly one of $\eta,\zeta$ is a unit in this residual
subcase. The exact stabilizer (AS5) exchanges these two alternatives.
It keeps $k$ a bounded unit, although a second orientation need not
keep $k\le1$. Choose the representative $\eta$ a unit and
$\zeta$ infinitesimal. Then
\[
 m/d,\beta/d\text{ are unbounded},\qquad
 C/d,A/d\text{ are strong},\qquad \alpha/d\text{ is weak}.
 \tag{NI7}
\]
Indeed $\theta=t\zeta$ is infinitesimal, $E$ is a unit,
$P_d,P_4,P_\alpha,r+\eta$ are units, and (NC4) gives
$m/d\asymp1/(M\zeta)$ and $\beta/d+1/2\asymp1/\theta$.
The reflected representative also has a unit $E$: its exact value
is $E'=h(1+y)/G$ in the original coordinates, with $h$ a unit by
the same $\theta=1$ exclusion used in the ordinary-angle proof.
There is no conclusion about the size of $D/d$ in (NI7).

If $v_c$ is unbounded, both $\eta,\zeta$ are noninfinitesimal and
at least one is unbounded. To see the first assertion, an infinitesimal
$\zeta$ would force $h\sim1$ and $\eta$ bounded by $E>0$;
an infinitesimal $\eta$ would force $r\sim1$ and $\zeta$ bounded
by $Q>0$. Two rate regimes are then useful.
\begin{itemize}
\item If $\eta$ is unbounded and $\zeta$ a unit, then $h$ is
infinitesimal. The exact alternative numerator identity
\[
 B_\beta=(x\eta h)\{K(1+\zeta)-x\theta^2\}
             +xrGT-\theta(1+2xh+x\zeta)             \tag{NI8}
\]
shows that $x\eta h$ is a unit: otherwise its residue at
$\theta=1$ would be $-1-x\zeta(1-r)<0$.
Thus $P_\alpha$ is a unit and $P_d\asymp\eta$.
Equations (NC4) give
\[
 C/d\asymp(1-r)^{-1},\quad A/d\asymp r^{-1},\quad
 m/d\asymp M^{-1},\quad\alpha/d\asymp Q^{-1},
 \quad D/d,\beta/d\text{ weak}.                    \tag{NI9}
\]
For the last assertion use $J\sim1$, $\theta\sim1$ and
$P_d/(Ky)\sim1$ in the original $D/d+1/2,\beta/d+1/2$
formulas. The total-scale condition $\delta_{\rm tot}$ infinitesimal
is therefore equivalent on this typed branch to at least one of
$r,1-r,M,Q$ being infinitesimal. If all four are units, the inherited
positive-total-scale theorem already applies.

\item If $\zeta$ is unbounded, $r$ is infinitesimal and $M$ a
unit. The noninfinitesimal $\eta$ obeys $\eta h\asymp1$.
For if $\theta$ is infinitesimal this follows from $E>0$ and
$h\sim1$; otherwise divide (NI8) by $\zeta$, use
$xrG/\zeta=O(r)$, and obtain that $x\eta h$ has a positive-unit
lower bound. Therefore (NC2) and (NC4) give
\[
 C/d\asymp\zeta^{-1},\quad m/d\asymp\zeta^{-1},\quad
 A/d\asymp r^{-1},\quad\alpha/d\asymp\zeta/Q.
 \tag{NI10}
\]
Thus $C,m$ are weak and $A,\alpha$ unbounded. No size restriction
on $D,\beta$ is inferred. This branch already has infinitesimal
$\delta_{\rm tot}$.
\end{itemize}
All implications retain the original physicality and exact source openings.


\section{\texorpdfstring{Exact C2 source and its uniqueness}{Exact C2 source and its uniqueness}}
\label{component:p6e5d41ed80b3}
\label{p6e5d41ed80b3:app:c2-source}
In this appendix, $c$ denotes the actual half-angle and
$C$ the event weight previously denoted $c$. Set
\[
 z=c^2,\quad A_c=1+zr,\quad
 h=b-a^2t,\quad g=(1+z)b-(1-r+a^2zt),\quad j=1+zr-t,
\]
\begin{align*}
\Delta={}&(1+z)zb^2+2z(1+z)br+(1+z)(1+a^2z)b\\
 &+z^2(1-a^2)r^2+2z(1+a^2z)r+1+a^2z,\\
N={}&(1+z)(b+r)(b+a^2(1-t))-a^2(1+a^2z)(1-r)t.
\end{align*}
The unchanged $d=1$, $\lambda=\beta=0$ source is
\begin{align*}
n&=\frac{(b+r)A_c(A_c+a^2zt)-r\Delta}{2r\Delta},&
m&=\frac{(1+a^2z)(1-r)A_c(A_c+a^2zt)}{2g\Delta},\\
C&=\frac{A_cN}{2a^2t\Delta},&
A&=\frac{a^2A_c(b+r)(A_c+a^2zt)}{2h\Delta},\\
D&=\frac{(1+z)(b+r)(b-a^2zr)(A_c+a^2zt)}{2a^2j\Delta},&
\alpha&=\frac{z^2(b+r)N}{2(1-zh)\Delta}.
\end{align*}
Its exact contact is
\[
 (u_c,v_c;W_c,X_c,Y_c,Z_c)
   =\left(c,ac;\frac ba,-cr,c\left(\frac ba-at\right),1\right).
\]
The positive chart requires $a,c,b,t>0$, $0<r<1$, $h,g,j>0$,
$1-zh>0$, $b-a^2zr>0$, and nonnegative original event weights.
The source identities and the full independent-test radial certificates
are included in the accompanying inherited modules; the source-entry
proof above imposes no hidden shape-unit hypothesis.


The numerator has the useful positive identity
\[
N=(1+c^2)(b+r)h+a^2\{g+j+rt(1+a^2c^2)\}>0.
\]
It follows by substituting the three source-gap definitions; every term
is positive in the domain used by the contact-entry theorem.
For the determinant calculation use the aliases
\[
 R_2=R_m,\quad R_3=uW-Y,\quad R_5=R_D,\quad R_\alpha=vY-Z,
 \qquad(e_1,e_2,e_3,e_4,e_5)=(n,m,C,A,D).
\]
\par\noindent\textit{Source inclusion. }The original input is reproduced once as Component p61209c7390b0 (Section~\ref{component:p61209c7390b0}).\par


\section{\texorpdfstring{A short determinant check for exact source entry}{A short determinant check for exact source entry}}
\label{component:p61209c7390b0}

\label{p61209c7390b0:sec:c2-contact-determinant}

Here is a reproducible uniqueness check for the six source coefficients.
Let $M$ have rows indexed by $(W,X,Y,Z,u,v)$ and columns equal to the
derivatives of
\[
D_uD_vX^2,\quad R_2^2,\quad D_vR_3^2,\quad
D_uD_vY^2,\quad R_5^2,\quad D_uR_\alpha^2,
\]
evaluated at $(u,v;W,X,Y,Z)=(c,ac;b/a,-cr,c(b/a-at),1)$.
Write $z=c^2$, $D=1+z$, $E=1+a^2z$, $h=b-a^2t$,
$g=Db-(1-r+a^2zt)$, and $j=1+zr-t$.
Factoring the six columns by
\[
2cr,\quad2cg,\quad2act,\quad-2ch/a,\quad2acj,
\quad2(1-zh)
\]
and the first row by $c$ leaves the matrix
\[
B=\begin{pmatrix}
0&a&E&0&-1&0\\
-DE&-1&0&0&-az&0\\
0&az&-E&-DE&1&-acD\\
0&-c&0&0&ac&D\\
zrE&zh-1&(b/a)E&-zhE/a&azr-b/a&c(1-zh)\\
azrD&b/a+zh/a&a^2zt&-zDh&1+zr&-cDh/a
\end{pmatrix}.
\]
Use the pivots in row $2$, column $1$; row $1$, column $3$; and
row $4$, column $6$.  Eliminate those three columns, multiply the
remaining second and third rows by $D$ and $E$, and factor $-DE$
from their middle column.  The determinant reduces to
$\det B=D^2E^2\det T$, where
\[
T=\begin{pmatrix}
a&1&a^2z\\
-1-b-z(r+a^2t)&zh/a&az(r-1+zh)\\
b/a+az(h-r)&zh&E(1+zh)+z(r+a^2t)
\end{pmatrix}.
\]
Subtract $(zh/a)$ times the first row from the second and $zh$ times
the first row from the third.  Expanding along the middle column gives
\[
\det T=[1+b+z(b+r)][1+z(a^2+b+r)]
-z(1-r)(b-a^2zr)=\Delta.
\]
The last equality follows by expanding the two products and is exactly
the source definition of $\Delta$.  Restoring the factors proves
\begin{equation}
\det M=-64a c^6rgt h j(1-zh)D^2E^2\Delta.
\label{p61209c7390b0:eq:c2-entry-determinant}
\end{equation}
Each source coefficient occurs with a factor $4$ in the parent, so the
stationary linear system in $(e_1,e_2,e_3,e_4,e_5,\alpha)$ has determinant
$4^6\det M$.  On the positive source domain this is nonzero.  Since the
six source formulas above satisfy all six stationary equations,
they are their unique solution.  This checks exact coefficient entry,
independently of the subsequent asymptotic Gram argument.


\section{\texorpdfstring{Self-contained C1 source definitions and contact verification}{Self-contained C1 source definitions and contact verification}}
\label{component:p0acdb66bba2f}

\label{p0acdb66bba2f:app:c1-source-definitions}

The letter $d$ in the contact--conductor proof denotes the new parameter
$\delta$, not the Bell coefficient. Here the latter is denoted $d_{\rm B}$,
and the original event coefficient $e_5$ is denoted $D_{\rm e}$.
The testing factor is always $D=1+u^2$. The conductor monomial
$C=\delta b^2uw$ is distinct from the clearing multiplier
$\mathcal C=K^6Q^4T_n$.

\label{component:pfecc204a8025}
\label{pfecc204a8025:app:c1-native}
\paragraph{Scope of the polynomial identities.} The seven coquotient formulas and the contact identities (SD9)--(SD11) are exact for arbitrary positive native parameters. The size assumptions below are needed only for the displayed bounded-Gram application.
\paragraph{Domain and seven native coquotients.}
The coefficient field and testing field are independent real closed
fields. On this C1 source branch, $k,f,g>0$ are coefficient
infinitesimals, $e=kf/g$ is positive unbounded, and
$\Delta=e^{-1}$. The shapes $a,t$ are positive bounded elements;
\[
 R=1+k^2(1-af^2),\qquad T=1+k^2-k^2g^2t,
 \qquad M=R-t>0,\quad J=T-a>0.
\]
All original event coefficients are assumed nonnegative. In particular,
the formulas below retain the requirements $D_{\rm e},\beta\ge0$.
No relation compares a testing rate with a coefficient rate.
Put
\[
 \kappa=k^2,\quad \eta=f^2,\quad \zeta=g^2,
 \quad V_c=\kappa\eta\zeta,\quad
 L_c=1-a\eta,\quad H_c=1+(1-a)\eta,\quad Q_c=1-\kappa\zeta+aV_c.
\]
Define
\begin{align*}
P_4={}&1+\kappa+\eta\kappa(a-\kappa-1)-\zeta\kappa t
 +V_c[-\zeta\kappa t^2+2\kappa t+a(1-t)+t],\\
P_\alpha={}&1+\kappa-a\eta\kappa
 +\zeta[(1+\kappa)-(2+\kappa)t]
 +V_c[-\zeta t^2+a(t-1)+t],\\
P_d={}&(1+\kappa)^2-a\eta\kappa^2
 +\zeta\kappa[(\kappa+1)-(\kappa+2)t]\\
 &+V_c[-\zeta\kappa t^2+\kappa a(t-1)+\kappa t+\kappa+1].
\end{align*}
The complete seven coquotients are
\begin{align*}
c_2&=2(1+\kappa)(1-t\zeta)L_c(1+V_c)H_cQ_c,\\
c_b&=2(1+\kappa)L_c[1+(1-t)\zeta]H_cQ_c,\\
c_3&=2(1+\kappa)L_cH_cQ_c(1-\kappa\eta+tV_c),\\
c_5&=2(1+\kappa)L_cH_cQ_c(1+V_c),\\
c_4&=2Q_cP_4,\qquad c_\alpha=2L_cH_cP_\alpha,
 \qquad c_d=4L_cQ_cP_d.
\end{align*}
Their residues, in this order, are $2,2,2,2,2,2,4$.
Here $c_b$ is the original $c_\beta$, and the subscript on $c_d$ is a
name rather than the conductor parameter.

\paragraph{Original source and full parent.}
The original event coefficients and multiplier are
\[
\begin{gathered}
\lambda=n=0,\qquad c=2,\qquad
 d_{\rm B}=\frac{2a\eta c_d}{c_3},\qquad
 m=\frac{2ae^2c_2}{Mc_3},\qquad A=\frac{2a\eta^2c_4}{c_3},\\
\alpha=\frac{2a\Delta^2k^6\eta^2c_\alpha}{c_3},\qquad
D_{\rm e}+\frac{d_{\rm B}}2=\frac{2ac_5}{Jc_3},\qquad
\beta+\frac{d_{\rm B}}2=\frac{2ae^2c_b}{tc_3},\qquad
\Lambda=\frac{RTc_3}{2a}.
\end{gathered}                                                    \tag{SD1}
\]
Let $F$ denote the original complete parent, including its signed core.
For a testing state $(W,X,Y,Z)$ define
\[
\begin{gathered}
 D=1+u^2,\quad D_v=1+v^2,\quad B=X+uZ,\quad Q_A=W+uY,\\
 r=uW-Y,\quad r_2=vQ_A-B,\quad r_5=r-v(Z-uX).
\end{gathered}
\]
With $H_2=4Tc_2$, $H_b=4Tc_b$, $H_3=4Rc_3$, $H_5=4Rc_5$,
the exact entire parent is
\begin{align*}
\Phi:=\Delta^2\Lambda F={}&H_2\frac RM r_2^2
 +H_b\frac RtD_v B^2+\Delta^2H_3\frac TaD_vr^2
 +\Delta^2H_5\frac TJr_5^2\\
 &+4RT\Delta^2\eta^2c_4DD_vY^2
 +4RT\Delta^4k^6\eta^2c_\alpha D(vY-Z)^2\\
 &+2RT\Delta^2\eta c_dD\{(vY-Z)^2-D_vZ^2\}.
                                                               \tag{SD2}
\end{align*}
Thus every original source event is accounted for.

\paragraph{Retained packets and state coordinates.}
In the remaining chart $Z=1$ with $u,v,W,X,Y$ infinitesimal, put
\[
 p=t/R,\quad h=a/T,\quad q_A=v(W+uY),\quad q_B=v(1-uX),
 \quad N=B-pq_A,\quad M_B=r-hq_B.
\]
Retain the four nonnegative packets
\[
 \Pi_1=H_2\frac tM N^2,\quad
 \Pi_2=H_b\frac MtD_vN^2,\quad
 \Pi_3=\Delta^2H_3\frac JaD_vM_B^2,\quad
 \Pi_4=\Delta^2H_5\frac aJM_B^2.                         \tag{SD3}
\]
Their subtraction is the exact identity
\begin{align*}
G=\Phi-\sum_j\Pi_j={}&H_2\{N^2-2q_AN+(1-p)q_A^2\}
 +H_bD_v\{N^2+2q_AN+pq_A^2\}\\
 &+\Delta^2H_3D_v\{M_B^2+2q_BM_B+hq_B^2\}
 +\Delta^2H_5\{M_B^2-2q_BM_B+(1-h)q_B^2\}
 +\mathcal W,\\
\mathcal W={}&4RT\Delta^2\eta^2c_4DD_vY^2
 +4RT\Delta^4k^6\eta^2c_\alpha D(vY-1)^2\\
 &+2RT\Delta^2\eta c_dD\{(vY-1)^2-D_v\}.
                                                               \tag{SD4}
\end{align*}
The testing coordinate $w=(W+uY)/D$ has the full inverse
\[
\begin{gathered}
 X=N+pDvw-u,\qquad q_B=v(D-uN-puDvw),\\
 W=w+\frac{u(M_B+hq_B)}D,\qquad
 Y=uw-\frac{M_B+hq_B}D.
\end{gathered}                                \tag{SD5}
\]
Its only testing denominator is the positive unit $D$.
On the negative-core branch, sign gauges give $u,v>0$, and $w>0$:
after eliminating the two positive core squares the remaining scalar is
\[
 \frac{uvDD_v}{\mathcal D}\{Duv w^2-2DD_vw+uv\},\qquad
 \mathcal D=u^2v^4+u^2v^2+u^2+v^4+2v^2+1>0,
\]
which is nonnegative if $w\le0$. The eliminated matrix is
$\left(\begin{smallmatrix}1+v^2+u^2v^2&uv\\uv&(D+v^2)/D\end{smallmatrix}\right)\succ0$.
The nonnegative-core branch and literal axes already give direct SOS.
Thus the radicals in the conductor proof separate between the two fields.

\paragraph{Coefficient substitution and complete first normals.}
Put
\[
 K=1+k^2,\quad \delta=\Delta R/K,\quad b^2=\eta K/R,
 \quad Q=K+ak^2b^2,\quad T_n=K^2-t\delta^2k^4Qb^2.
\]
Then $\delta,k,b$ are positive infinitesimals, and
\[
 \eta=Kb^2/Q,\quad \Delta=\delta Q/K,\quad
 \zeta=\delta^2k^2Qb^2/K,\quad R=K^2/Q,\quad T=T_n/K.     \tag{SD6}
\]
All displayed denominators are positive pure-coefficient units.
Set $C_i=K^2Q^3c_i$ after this substitution. Explicitly, a native monomial
$c\,k^jf^{2r}g^{2s}a^\alpha t^\tau$ becomes
\[
 c\,\delta^{2s}k^{j+2s}b^{2(r+s)}a^\alpha t^\tau
       K^{2+r-s}Q^{3-r+s}.                            \tag{SD7}
\]
The native coquotients have $-2\le r-s\le3$, so these are polynomials
with no negative exponents or shape inverses.
Define $V_n,Y_n,B_0,A_c$ as in the conductor proof and
$\mathcal C=K^6Q^4T_n$. The factored scalar printed there is exactly
$F_0=(\mathcal CG)|_{N=M_B=0}$. For completeness put
\[
 S_n=8\delta^2K^2b^4C_4D_vY_n
 +8\delta^4Q^2k^6b^4C_\alpha vB_0
 +4\delta^2KQb^2C_dvB_0.
\]
The complete coefficients on the hard normals $N,\delta M_B$ are
\begin{align*}
 L_1={}&8K^3QT_n^2(C_bD_v-C_2)Dvw
       -8\delta^2K^2Q^2uD v^2V_nA_c+K^3auvS_n,\\
 L_2={}&8\delta K^2Q^2T_nD vV_n(C_3D_v-C_5)-K^2T_nS_n/\delta.
                                                               \tag{SD8}
\end{align*}
The quotient in the last line is polynomial because $S_n$ contains
$\delta^2$. These formulas follow by substituting (SD5) into (SD4),
channel by channel; they include all normal terms of degree one.

\paragraph{Short native proof of the exact contact identities.}
Write $V=1+V_c$ and $K=1+\kappa$ in this paragraph. The factored
coquotients above give immediately
\[
 c_b-c_2=\zeta c_3,\qquad c_3-c_5=-\kappa\eta c_b,
 \qquad c_2=(1-t\zeta)c_5.                            \tag{SD9}
\]
Hence $Vc_b-c_2=\zeta c_5$ and
$Vc_3-c_5=-\kappa\eta c_2$.
The remaining three balances are the following short polynomial
identities; substituting the displayed $P_4,P_\alpha,P_d$ proves them
by multiplication and collecting terms:
\begin{align*}
\mathfrak Y:={}&2\kappa c_2+2VL_cc_4
                 -2\kappa^2\zeta Q_cc_\alpha-Q_cc_d=0,\\
\mathfrak S:={}&\kappa L_cc_2+c_5+VL_c^2c_4+\kappa Q_c^2c_\alpha
                 +\tfrac12(\kappa\zeta L_c^2-2L_c-\eta)c_d=0,\\
\mathfrak V:={}&\kappa Rc_2+t\kappa V_cc_b+aV_cc_3+Tc_5
 +KV_cL_c^2c_4-\kappa^2\zeta KQ_cL_cc_\alpha
 -\tfrac K2(Q_cL_c+\eta)c_d=0.                        \tag{SD10}
\end{align*}
For example, after the common factor $4L_cQ_c$ is removed,
the first balance is just
\[
 \kappa K(1-t\zeta)VH_c+VP_4
       -\kappa^2\zeta H_cP_\alpha-Q_cP_d=0.
\]
The two other balances are equally direct identities in the same three
short polynomials; no expanded scalar or ideal calculation enters.

The exact contact, in the original coordinates of (SD2), is
\[
 u=k,\quad v=\Delta k^2\eta,\qquad
 (W,X,Y,Z)=(\Delta k,-k(1-t\zeta),\Delta k^2L_c,1).      \tag{SD11}
\]
At this point
\[
 r_2=Mk\zeta,\quad B=tk\zeta,\quad r=av,\quad
 r_5=-Jv,\quad vY-1=-Q_c,\quad D=K,\quad D_v=V.
\]
Set $\Psi=\Phi/(4RT)$ and let $E_s=\tfrac12\partial_s\Psi$,
with all other original testing coordinates held fixed.
Direct differentiation of the seven terms in (SD2), followed by (SD9),
gives at (SD11)
\begin{align*}
 \Psi&=K\zeta^2\mathfrak S,\\
 E_W&=kv\{\zeta c_2+\Delta^2(Vc_3-c_5)\}=0,\\
 E_X&=k\zeta\{Vc_b-c_2-\zeta c_5\}=0,\\
 E_Y&=\tfrac12K\Delta^3\kappa\eta^2\mathfrak Y,\\
 E_u&=k\zeta^2\{\mathfrak S-c_2+(1-t\zeta)c_5\},\\
 E_v&=v\Delta^2\mathfrak V.
\end{align*}
They all vanish by (SD9)--(SD10). Homogeneity in $(W,X,Y,Z)$ also gives
$E_Z=0$. This proves contact and full first-order stationarity directly
in the native source.

In (SD5)--(SD6), this point is exactly
\[
 (u,w,v,N,M_B)=(k,\delta k,\delta k^2b^2,0,0).
\]
All four packets (SD3) and their first derivatives vanish there.
Changing coordinates and multiplying by $\mathcal C$ therefore gives
$(F_0)_c=(\partial_uF_0)_c=(\partial_wF_0)_c
=(\partial_vF_0)_c=(L_1)_c=(L_2)_c=0$.
The leading polynomial $P$ is a sum of circuit squares vanishing at
this contact, so $P\in J^2$. This proves every identity (CC7) without
using the expanded-source checker as a premise.

\paragraph{Positive correction weight without a large expansion.}
Use the weights $(2,2,6,0,0,2,4,18)$ on
$(\delta,k,b,a,t,u,w,v)$. Write $O_{\rm wt}(s)$ for a polynomial all
of whose monomials have weight at least $s$. Native $\eta$ and $\zeta$
have weights at least $12$ and $20$. Setting them to zero in the seven
short coquotients, and then applying (SD7), gives
\[
 C_i=2K^6+O_{\rm wt}(12)\quad(i=2,b,3,5,4,\alpha),\qquad
 C_d=4K^7+O_{\rm wt}(12).
\]
Also
\[
 Q=K+O_{\rm wt}(16),\quad T_n=K^2+O_{\rm wt}(24),\quad
 V_n=K^2+O_{\rm wt}(24),\quad D_v=1+O_{\rm wt}(36),
\]
\[
 Y_n=O_{\rm wt}(6),\quad B_0=-K^3+O_{\rm wt}(24),\quad
 A_c=2K^8+O_{\rm wt}(12),\quad
 \mathcal C=K^{12}+O_{\rm wt}(16).
\]
In the five factored scalar terms, discard only coefficients carrying
weight at least $12$. The surviving expression is exactly
\[
 8K^{14}D\{D(A^2+B^2)+K(C^2+D_0^2)-2K^2BC\}
       =8K^{14}DP.
\]
Here $A^2,D_0^2$ have weight $44$, while $B^2,C^2,BC$ have weight $40$.
Consequently $E_0=O_{\rm wt}(52)$, sufficient for every scalar lift
coefficient to have weight at least $8$.
For the normals use (SD9) in its polynomial forms
\[
 K(C_b-C_2)=\delta^2k^2Qb^2C_3,\qquad
 Q(C_3-C_5)=-k^2Kb^2C_b.
\]
Since $K,Q$ have constant term one, these show weights at least $20,16$
for the two differences, respectively, without putting their inverses
in any lifting coefficient. The displayed $S_B,S_C$ in the conductor
proof have weight at least $12$. Its factored normal formulas now give
$L_1=O_{\rm wt}(42)$ and $L_2=O_{\rm wt}(32)$.
The monic lift therefore has positive coefficient weights at least
$8,20,10$. The stronger exact checker bounds $12,20,12$ are optional.

\paragraph{Full hard block and restoration.}
There are no state-normal terms above degree two. Put
\[
 C_B=H_3D_vh+H_5(1-h),\qquad
 B_Y=4RT\Delta^2\eta^2c_4DD_v
 +4RT\Delta^4k^6\eta^2c_\alpha Dv^2
 +2RT\Delta^2\eta c_dDv^2.
\]
Before multiplication by $\mathcal C$, the exact quadratic block on
$(N,\delta M_B)$ has entries
\begin{align*}
 A_{11}&=H_2+H_bD_v+\Delta^2C_Bu^2v^2+B_Yh^2u^2v^2/D^2,\\
 A_{12}&=-\frac{\Delta^2}{\delta}(H_3D_v-H_5)uv
                  -\frac{B_Y}{\delta D^2}huv,\\
 A_{22}&=\frac{\Delta^2}{\delta^2}(H_3D_v+H_5)
                  +\frac{B_Y}{\delta^2D^2}.
\end{align*}
The displayed ratios cancel coefficient monomials exactly:
$\Delta/\delta=Q/K$ is a positive unit and
$B_Y/\delta^2$ is bounded infinitesimal. Thus this full block has
residue $16I_2$, also after multiplication by $\mathcal C$.
The conductor proof and the positive circuit frame give the bounded
strict-residue Gram for $\mathcal CG$. Restore the four full positive
packets $\mathcal C\Pi_j$ from (SD3), then undo
$\mathcal C\Delta^2\Lambda$. This recovers the entire original parent;
retained gap ratios need not be bounded and are never placed in that
bounded Gram. The complementary state and outer-angle charts require
the already proved C1 chart arguments. This appendix makes no claim of
an exhaustive global source cover.


\section{\texorpdfstring{A short contact--conductor proof of the C1 memberships}{A short contact--conductor proof of the C1 memberships}}
\label{component:pfabc29e28ab3}
% Optional module: the source notation is that of (UZ4)--(UZ6).
% All statements below are exact; no localized ring is used.


Write $d=\delta$ in this section; this local parameter is distinct from
the global Bell coefficient $d_{\rm B}$. Put
\[
 R_0=\mathbb Q[d,k,b,a,t],\qquad R=R_0[u,w,v],\qquad
 J=(u-k,w-dk,v-dk^2b^2).
\]
For $F\in R$, write $F_c=F(k,dk,dk^2b^2)$. Define
\[
 A=wv,\quad B=dv,\quad C=db^2uw,\quad D_0=d^2k^3b^2,
 \qquad M=(A,B,C),
\]
\[
 x=B-C,\qquad y=A-kC,\qquad z=uB-kC,\qquad h=kC-D_0.
\]
These are the original roots $g_0,h_{13},h_{23},h_{34}$; in particular,
$D_0=a_4^2$, whereas the source monomial in $g_4$ is $kD_0$.
Adjoin the separated radicals used in the original four-monomial circuit,
\[
 a_1=\sqrt{wv},\quad a_2=\sqrt{duv},\quad
 a_3=b\sqrt{dkuw},\quad a_4=dk^{3/2}b,
 \qquad q=a_1a_2-a_3a_4.
\]
Only $q^2\in R$ is needed in the scalar lift. The following three identities
are the entire conductor calculation:
\begin{align}
 C(u-k)&=z-ux,\label{pfabc29e28ab3:CC-u}\\
 C(w-dk)&=dy-wx,\label{pfabc29e28ab3:CC-w}\\
 C(v-dk^2b^2)&=b^2(wz+dky-kwx).\label{pfabc29e28ab3:CC-v}
\end{align}
Thus $CJ\subset(x,y,z)$ without inverting any parameter.
For
\[
 r_*=y+z-h-2kx
\]
we also have the two exact identities
\begin{align}
 r_*&=(v-dk^2b^2)\bigl((w-dk)+d(u-k)\bigr)
       -dkb^2(u-k)(w-dk),\label{pfabc29e28ab3:CC-star}\\
 kCr_*&=q^2-yz.\label{pfabc29e28ab3:CC-repair}
\end{align}
The first says $r_*\in J^2$. The second supplies the otherwise missing
conductor factor through one of the original circuit squares.

\paragraph{Explicit monic division.}
Every $S\in R$ has a polynomial decomposition
\[
 S=S_c+(u-k)T_u+(w-dk)T_w+(v-dk^2b^2)T_v.                 \tag{CC1}
\]
One obtains it by substituting $u=k$, then $w=dk$, then $v=dk^2b^2$,
and using $X^n-Y^n=(X-Y)\sum_{j=0}^{n-1}X^{n-1-j}Y^j$ at each step.
There are no parameter denominators. If $S_c=0$, equations
\eqref{pfabc29e28ab3:CC-u}--\eqref{pfabc29e28ab3:CC-v} give
\[
 CS=x(-uT_u-wT_w-b^2kwT_v)
     +y(dT_w+dkb^2T_v)+z(T_u+b^2wT_v).                 \tag{CC2}
\]

\begin{lemma}[First-order and restricted second-order contact lifts]
In the above ring, the following assertions hold, with polynomial
coefficients in $R$:
\[
 M\cap J\subset (x,y,z),\qquad
 \bigl(M^2+(D_0^2)\bigr)\cap J^2
       \subset (x,y,z,h,q)^2.                         \tag{CC3}
\]
Here $(D_0^2)$ is the principal ideal generated by $D_0^2$; the second
assertion does not replace it by all mixed products involving $D_0$.
\end{lemma}
\begin{proof}
For the first assertion, write $F=R_AA+R_BB+R_CC$. Since
$A=y+kC$ and $B=x+C$,
\[
 F=R_Ay+R_Bx+CS,\qquad S=kR_A+R_B+R_C.
\]
The contact value is $F_c=C_cS_c$, where $C_c=d^2k^2b^2$ is a nonzero
polynomial. Since $R_0$ is a domain, $F_c=0$ implies $S_c=0$.
Apply (CC2). This cancellation proves a zero identity; it introduces no
coefficient $1/C_c$.

For the second assertion choose any displayed representation
$F=\sum_{i\le j}f_{ij}m_im_j+\rho D_0^2$, with
$(m_1,m_2,m_3)=(A,B,C)$. Substitute
\[
 A=y+kC,\qquad B=x+C,\qquad D_0=kC-h.
\]
This gives
\[
 F=Q_2+C(xH_x+yH_y+zH_z+hH_h)+C^2S,\qquad H_h=-2k\rho,       \tag{CC4}
\]
where $Q_2\in(x,y,h)^2$ and initially $H_z=0$.
The contact value again implies $S_c=0$. Use (CC2) on $CS$ and absorb
$C^2S$ into the four coefficients in (CC4). This leaves $H_h=-2k\rho$
unchanged and gives the same formula without $C^2S$.

The first derivatives of $F\in J^2$ vanish at contact. The first three
root gradients have determinant
\[
 \det\left(\frac{\partial(x,y,z)}{\partial(u,w,v)}\right)_c
       =-d^4k^4b^4\ne0.                              \tag{CC5}
\]
Equation~\eqref{pfabc29e28ab3:CC-star} gives a gradient relation with coefficient vector
$(-2k,1,1,-1)$. Its coefficient at $h$ is $-1$.
Set $\lambda=2k\rho_c$. Subtract
$\lambda(-2k,1,1,-1)$ from $(H_x,H_y,H_z,H_h)$.
The last resulting contact value is zero, and the first three satisfy the
homogeneous linear system with determinant (CC5). Multiplication by its
adjugate and cancellation in the domain $R_0$ show that those three
contact values are zero too. No inverse of the determinant is used in a
coefficient: each resulting polynomial $J_i$ lies in $J$ and is divided
only by the monic procedure (CC1).
Consequently
\[
 F=Q_2+2k\rho_cCr_*+C\sum_{i=1}^4r_iJ_i
   =Q_2+2\rho_c(q^2-yz)+\sum_{i=1}^4r_i(CJ_i),          \tag{CC6}
\]
where $(r_1,r_2,r_3,r_4)=(x,y,z,h)$. Each $CJ_i$ is explicitly a
combination of $x,y,z$ by (CC2), which proves the assertion in the
unlocalized polynomial ring.
\end{proof}

\paragraph{Why the restriction matters.}
The unrestricted claim obtained by replacing $M^2+(D_0^2)$ with
$(A,B,C,D_0)^2$ is false. Indeed $Cr_*\in(A,B,C,D_0)^2\cap J^2$.
Under the positive weights below it has weight $42$. Every weight-$42$
member of the square of the original soft-root ideal must be divisible
by $x$, since $x$ has weight $20$ and every other original soft root has
weight at least $22$. But $Cr_*$ is not divisible by $x$: setting
$d=k=b=w=1$, $u=2$, $v=2$ gives $x=0$ and $Cr_*=2$.
The factor $k$ supplied by $D_0^2$ in (CC4) is therefore substantive.

\paragraph{Application to the complete source scalar and both normals.}
Use precisely the seven source coquotients and the definitions in
Appendix~\ref{p0acdb66bba2f:app:c1-source-definitions}, renaming their clearing multiplier $\mathcal C=K^6Q^4T_n$
to distinguish it from the present conductor monomial $C$. In particular
\[
 K=1+k^2,\quad Q=K+ak^2b^2,\quad T_n=K^2-td^2k^4Qb^2,
 \quad D=1+u^2,\quad D_v=1+v^2,
\]
\[
 V_n=K^2-tQuvw,\quad Y_n=KT_nuw-avV_n,
 \quad B_0=vY_n-KT_n,
 \quad A_c=C_5T_n+aK(C_3D_v-C_5).
\]
The exact scalar $F_0$ has the following five-term form:
\begin{samepage}
\begin{align*}
F_0={}&4KQT_n^2[C_2(K^2-tQ)+C_bD_vtQ]D^2A^2\\
 &+4Q^2A_cD^2V_n^2B^2\\
 &+4K^3C_4DD_v(KT_nC-ab^2BV_n)^2\\
 &+4KQ^2C_aDB_0^2D_0^2\\
 &+2K^2QC_dD\{b^2B^2(Y_n^2-K^2T_n^2+2aKT_nV_n)
                         -2K^2T_n^2BC\}.
\end{align*}
\end{samepage}
Also the exact leading polynomial is
\[
 P=D(A^2+B^2)+K(C^2+D_0^2)-2K^2BC.
\]
Thus $E_0=F_0-8K^2D(K^6Q^4T_n)P$ belongs to
$M^2+(D_0^2)$ directly, without ideal division.

For the first normals put
\begin{align*}
 S_B={}&-8aK^2b^4C_4D_vV_n
          +8d^2Q^2k^6b^4C_aB_0+4KQb^2C_dB_0,\\
 S_C={}&8K^3T_nb^2C_4D_v.
\end{align*}
The normal numerator satisfies $S_n/d=BS_B+CS_C$. Hence
\begin{align*}
 L_1={}& A\,[8K^3QT_n^2(C_bD_v-C_2)D]\\
 &+B\,[-8dK^2Q^2uDvV_nA_c+dK^3auvS_B]
       +C\,[dK^3auvS_C],\\
 L_2={}&B\,[8K^2Q^2T_nDV_n(C_3D_v-C_5)-K^2T_nS_B]
       -C\,K^2T_nS_C.
\end{align*}
These formulas prove $L_1\in M$ and $L_2\in(B,C)$.
The native contact and derivative identities (SD9)--(SD11) in
Appendix~\ref{p0acdb66bba2f:app:c1-source-definitions} prove
\[
 (E_0)_c=(\partial_uE_0)_c=(\partial_wE_0)_c
          =(\partial_vE_0)_c=0,\qquad (L_1)_c=(L_2)_c=0.      \tag{CC7}
\]
These are source polynomial identities, with no positivity hypothesis or
parameter inversion. Their proof uses only six short native balances and
the chain rule. 
By (CC3),
\[
 E_0\in(x,y,z,h,q)^2,\qquad L_1,L_2\in(x,y,z).          \tag{CC8}
\]
Thus all three memberships previously expressed by the finite
Gr\"obner-basis witness have a single contact--conductor mechanism.

\paragraph{Bounded coefficients.}
Give $(d,k,b,a,t,u,w,v)$ weights $(2,2,6,0,0,2,4,18)$.
Then $A,y,z,h,q$ have weight $22$, while $B,C,x$ have weight $20$.
The three substitutions defining contact preserve weight, as do the
monic geometric-sum divisions after accounting for the weight of their
linear factor. The factored weight argument in
Appendix~\ref{p0acdb66bba2f:app:c1-source-definitions} gives, without a large expansion,
\[
 \operatorname{wt}(E_0)\ge52,\quad
 \operatorname{wt}(L_1)\ge42,\quad
 \operatorname{wt}(L_2)\ge32.
\]
Consequently the construction gives lower bounds $8,20,10$ for the
three families of correction coefficients, respectively. These bounds are positive. Every coefficient has nonnegative
exponents, including in the bounded shapes $a,t$; every correction
therefore has zero residue on the stated C1 branch. The full hard-normal
block and its positive residue are printed in
Appendix~\ref{p0acdb66bba2f:app:c1-source-definitions}. Restoration is exactly that of the source theorem: restore
$\mathcal C=K^6Q^4T_n$ times the four positive packets subtracted there,
then undo the positive multiplier $\mathcal C\Delta^2\Lambda$. The present argument proves the complete
memberships on that branch and does not assert a global source cover.


\section{\texorpdfstring{The C1 circuit needs six soft directions}{The C1 circuit needs six soft directions}}
\label{component:p58894e12bcef}

Let
\[
a_1=\sqrt{wv},\quad a_2=\sqrt{\delta uv},\quad
a_3=b\sqrt{\delta kuw},\quad a_4=\delta k^{3/2}b.
\]
The radicals are separated coefficient/testing radicals. Put
\[
h_{13}=a_1^2-a_3^2,\quad h_{23}=a_2^2-a_3^2,\quad
h_{34}=a_3^2-a_4^2,
\qquad q_1=a_1a_2-a_3a_4,\quad q_2=a_1a_3-a_2a_4.
\]
Averaging the three-square arithmetic--geometric-mean identities for
pairings $12|34$ and $13|24$ gives exactly
\begin{equation}
\sum_{i=1}^4a_i^4-4a_1a_2a_3a_4
=h_{13}^2+h_{23}^2+h_{34}^2-h_{13}h_{23}+h_{23}h_{34}
 +q_1^2+q_2^2.\label{p58894e12bcef:eq:compact-circuit}
\end{equation}
The contrast block is
\[
G=\begin{pmatrix}1&-1/2&0\\-1/2&1&1/2\\0&1/2&1\end{pmatrix}\succ0;
\qquad \det G_{[1]}=1,\quad\det G_{[2]}=3/4,\quad\det G=1/2.
\]
Unlike a single pairing, the two pairings control all three independent
contrasts of the four squared monomials.

Write $K=1+k^2$, $D=1+u^2$ and
\[
g_0=\delta(v-b^2uw),\qquad g_4=uwv-\delta^2k^4b^2.
\]
The leading scalar used in the complete source is
\[
P=D(w^2+\delta^2)v^2+\delta^2Kb^4u^2w^2+
\delta^4Kk^6b^4-2\delta^2K^2b^2uwv
=g_0^2+g_4^2+\sum_i a_i^4-4\prod_i a_i.
\]
There is a second exact simplification:
\begin{equation}
g_4=u h_{13}+k h_{23}-ku g_0+k h_{34}.\label{p58894e12bcef:eq:g4-conductor}
\end{equation}
All four displayed coefficients are infinitesimal in the allowed chart.
Thus $g_4^2$ can be retained as a positive square or absorbed as a
zero-residue Gram correction. Equations \eqref{p58894e12bcef:eq:compact-circuit} and
\eqref{p58894e12bcef:eq:g4-conductor} give a strict residue Gram on just
\[
\boxed{(g_0,h_{13},h_{23},h_{34},q_1,q_2).}
\]
The first-order and restricted second-order conductor theorem places
every source correction in the first five of these roots. The two hard
state normals therefore give an \emph{eight-root} bounded strict-residue
Gram for the normalized remainder after retaining the four positive
packets, replacing its previous thirteen-root frame. Restoring those
packets gives the entire parent. Their possibly unbounded gap ratios
are kept outside this bounded Gram. This is a local Gram-frame
dimension, not a global SDP block-size bound or a claimed optimal number
of tensor squares.

One pairing alone does not suffice. At the formal complex point
$(a_1,a_2,a_3,a_4)=(1,1,i,-i)$, its roots
$a_1^2-a_2^2$, $a_3^2-a_4^2$, $a_1a_2-a_3a_4$ all vanish, whereas
$h_{13}=2$. Hence their ideal does not contain $h_{13}$, and their squared
ideal cannot absorb every infinitesimal multiple of $h_{13}^2$.
This is an algebraic obstruction to that proposed reduction, not a physical
complex-state counterexample. The second pairing supplies exactly the
missing strict contrast direction.


\section{\texorpdfstring{A three-monomial contact lift with one unit coordinate}{A three-monomial contact lift with one unit coordinate}}
\label{component:p466b31015532}

\label{p466b31015532:sec:three-circuit}

Work in the same tensor ring $T=K_c\otimes_{\mathbb R}K_t$ and bounded
subring $B$ as throughout this package. Every inverse and radical below
is of a pure-field or positive separated element. Let $k>0$ be a
pure-source infinitesimal and $s>0$ a positive pure-source bounded unit.
Let $u>0$ be a pure-testing infinitesimal and $t>0$ a positive
pure-testing bounded element (possibly infinitesimal). Define
\[
 A_1=u^2t,\qquad A_2=\frac{ku t^2}{s},\qquad A_3=k^2s,
 \qquad G=kut,\qquad d_i=A_i-G.
\]
Then $A_1A_2A_3=G^3$. Two exact conductor identities are
\begin{equation}\label{p466b31015532:T3:conductors}
 G(u-k)=k d_1,\qquad G(t-s)=s d_2.
\end{equation}
They permit a unit testing ratio $t$ without comparing a mixed tensor.

Put $Q=G\sum_i d_i$ and $L=\sum_i A_i^2-3G^2$.
For $a_i=A_i^{1/3}$, the elementary cubic identity gives
\begin{align}\label{p466b31015532:T3:sos}
 Q&=\frac12\sum_{1\le i<j\le3}\sum_{\ell=1}^3
       \{\sqrt{G a_\ell}(a_i-a_j)\}^2,\\
 L&=\sum_i d_i^2+2Q.\nonumber
\end{align}
Thus $L$ has a positive constant diagonal Gram on a finite list of
three difference roots and nine cubic-identity roots. No assertion
that every nonnegative tensor is SOS is used.

\paragraph{Exact finite lift.}
Let $\lambda\in B_c$ have positive real residue, and consider
\[
 F=\sum_{i=1}^3(\lambda+\epsilon_i)A_i^2
       -(3\lambda+\epsilon_g)G^2,
 \qquad h=\sum_i\epsilon_i-\epsilon_g.
\]
Write $\Delta_u=u-k$, $\Delta_t=t-s$. Assume explicit finite identities
\begin{align}
 \epsilon_i-\epsilon_{i0}&=K_{iu}\Delta_u+K_{it}\Delta_t,
       &\epsilon_i,\epsilon_{i0}&\in B_{\rm Inf},\label{p466b31015532:T3:coeff}\\
 h&=b_u\Delta_u+b_t\Delta_t+
          \sum_{j,\ell\in\{u,t\}}H_{j\ell}\Delta_j\Delta_\ell,
       &b_u,b_t&\in K_c,\label{p466b31015532:T3:taylor}\\
 2(\epsilon_{10}-\epsilon_{30})+k b_u&=0,
 &2(\epsilon_{20}-\epsilon_{30})+s b_t&=0.\label{p466b31015532:T3:jets}
\end{align}
Require $K_{iu},H_{uu},H_{ut},H_{tu}\in B$ and
$K_{it},H_{tt}\in B_{\rm Inf}$. The distinction is essential:
the second conductor has a unit coefficient, so boundedness alone of
the second-coordinate arrays does not suffice.

Then $F$ is tensor SOS. Indeed, expanding $A_i=G+d_i$ and applying
\eqref{p466b31015532:T3:conductors} and \eqref{p466b31015532:T3:jets} yields the exact identity
\begin{align}\label{p466b31015532:T3:error}
 F-\lambda L={}&\sum_i\epsilon_i d_i^2+2\epsilon_{30}Q
 +2\sum_i(kK_{iu}d_1d_i+sK_{it}d_2d_i)\\
 &+k^2H_{uu}d_1^2+ks(H_{ut}+H_{tu})d_1d_2
          +s^2H_{tt}d_2^2.\nonumber
\end{align}
In detail, the two bare-$G$ first-order terms are cancelled by
$G^2b_u\Delta_u=kb_uGd_1$ and
$G^2b_t\Delta_t=sb_tGd_2$. All remaining coefficients in
\eqref{p466b31015532:T3:error} are bounded infinitesimals. The strict diagonal
Gram for $\lambda L$ in \eqref{p466b31015532:T3:sos} therefore remains strictly
positive in real residue. The bounded-Gram lemma supplies a finite
tensor SOS. This argument uses no power series or mixed inverse.

\paragraph{Normal extension.}
Suppose a coupling to a normalized hard state normal has the exact form
$\ell=\sum_i c_iA_i$ with $c_i\in B_{\rm Inf}$,
$\sum_i c_{i0}=0$, and
\[
 \sum_i(c_i-c_{i0})=J_u\Delta_u+J_t\Delta_t,
 \qquad J_u\in B,\quad J_t\in B_{\rm Inf}.
\]
Then
\[
 \ell=\sum_i c_i d_i+kJ_u d_1+sJ_t d_2
\]
has infinitesimal coefficients on the same root list. It can be joined
to a hard normal block with strictly positive real residue. Extra
already-controlled scalar roots may be appended, provided their
coupling coefficients are likewise infinitesimal. The source-specific
arrays and hard-block bounds remain separate proof obligations.

\paragraph{Small-$L$ coordinates.}
For $x=k^2$, $s=\sqrt{\eta\zeta}$ and $t=v/u$, this gives
\[
 (A_1,A_2,A_3,G)=
 \left(uv,\sqrt{\frac{x}{\eta\zeta}}\frac{v^2}{u},
       \sqrt{\eta\zeta}\,x,\sqrt{x}\,v\right).
\]
The lemma applies only when the stated pure-field size assumptions and
all finite scalar and normal identities have been proved. A leading
three-term AM--GM form alone does not close a source domain.


\section{\texorpdfstring{The source, the original polynomial, and the proof language}{The source, the original polynomial, and the proof language}}
\label{component:p472f434581f3}


Let $K_c,K_t$ be independent real closed extensions of $\mathbb R$.
Source parameters lie in $K_c$, and angles and state coordinates in
$K_t$. The assertion is a finite sum of squares in
$K_c\otimes_{\mathbb R}K_t$, for each such pair of configurations.
No order or general inverse is imposed on this tensor ring.
A bounded quantity has a finite real residue; a positive unit is
bounded above and below by positive reals. All asymptotic statements
about a source or a test are made inside its own field.

\paragraph{The original source.}
Keep the physical stationary C1 source, its strict gaps,
nonnegative event weights and all original residual intersections.
Normalize its positive coefficient $d$ to $1$. On C14.2,
$x=k^2>0$ is infinitesimal, $L=x\zeta$ is a positive bounded unit,
and $\eta$ is unbounded. Set
\[
 \sigma=\eta^{-1},\ z=L^{-1},\ b=Lr,\ Q=1-b,
 \theta=1-h,\ e=xh/\sigma,\ E=1-e,\ K=1+x.
\]
The source has $0<r,h,b<1$, $r=bz$, and its original strict gaps.
The following exact formulas are NC1--NC4 in these coordinates:
\[
\begin{split}
 J&=1+xh-\sigma(1-r),\qquad M=1+xr-xz\theta,\\
 P&=\sigma+K+\sigma z(1+bx+2xh)
                  +\sigma bx^2hz^2+x^2h\theta z,\\
 p=Ar&=\frac{z(1+xh)E+(1-r)(1+xhz)}{2P},\\
 l=\alpha zQ&=\frac{(1+\sigma r)
     \{e(1+xr)(1+\sigma z)+ME\}}{2P},\\
 C&=\frac{Kz(1+\sigma r)E}{2(1-r)P},\\
 D&=\frac{K(1+\sigma z)(1+\sigma r)-JP}{2JP},\\
 \beta\theta&=\frac{K(1+\sigma r)(1+xhz)-\theta P}{2P},\\
 \kappa&=\frac{K(1+\sigma r)(1+\sigma z)}{2MP},
 \qquad m=\kappa x^2z\frac h\sigma.
\end{split}                                                     \tag{P1}
\]
Thus $a=Az=p/b$ and $\ell=\alpha z=l/Q$; the symbol $p$ here
denotes a source product, not a testing angle. All original gaps
are positive, while $D,\beta$ may be zero.

\paragraph{The shared source sphere.}
Put $H_0=1+z-r=1+zQ$,
$B=P+xz(J-\sigma)$, $R_s=1+xr$, and
\[
 \mathcal H=R_sH_0+z\{(1-r)xh(2+xr)
        -e[2+x(1+h+r\theta)]+Ke^2\}.
\]
The exact source identities BE3--BE8 are
\[
\begin{gathered}
 2\theta P\beta=hB-\sigma H_0,\qquad
 2P(JD+x\theta\beta)=\sigma^2\mathcal H,\\
 d_0^2=\frac{2PJD}{\sigma^2\mathcal H},\qquad
 b_0^2=\frac{2Px\theta\beta}{\sigma^2\mathcal H},\qquad
 d_0,b_0\ge0,\quad d_0^2+b_0^2=1,\\
 H=\frac h\sigma=\frac{H_0}{B}
       +\frac\sigma x\frac{\mathcal H}{B}b_0^2.
\end{gathered}                                                   \tag{P2}
\]
They follow by clearing the positive source denominators in P1.
In this branch $e=O(x+\sigma)$ is infinitesimal, $M,P,J,B\sim1$,
$\mathcal H\sim H_0$, and $\kappa\sim1/2$. To obtain the only
non-immediate bound $e=O(x+\sigma)$, expand the original $D\ge0$
numerator in native $\eta=1/\sigma$ coordinates as in BE1:
\[
 e\{K+x\theta(1+xh)/L\}
       \le x(b_1/L)+\sigma(a_0/L),
\]
where
\[
\begin{split}
 b_1={}&-Lhrx-Lh-Lr+L-h^2rx^2+h^2rx-3h^2x\\
       &-2hrx+hx-3h+1,\\
 a_0={}&Kr+(1-r)\{1+L+2xh+rxL+rx^2h\}.
\end{split}
\]
Both right-hand source ratios are bounded and the left bracket is
at least one. P1 then gives the opening-uniform comparisons
\[
 C(1-r)/z\sim\tfrac12,\quad p\sim H_0/2,\quad l\sim\tfrac12,
 \qquad D=O(\sigma^2),\quad x\beta\theta=O(\sigma^2).
                                                               \tag{P3}
\]
In particular $C,A,\alpha$ have positive real lower bounds,
although their upper bounds may diverge. There is no inverse of a
sphere coordinate and no lower bound on $D$ or $\beta$.

\paragraph{The entire original polynomial.}
On the reciprocal Bob chart $w=1/v$, put $D_u=1+u^2$ and
$D_w=1+w^2$. For a real state $(W,X,Y,Z)$ define
\[
\begin{gathered}
 M_w=W-wX+uY-uwZ,\quad R_w=uW-Y,\\
 J_w=uwW+uX-wY-Z,\quad B_w=X+uZ,\\
 K_w=J_w^2+D_u(Y-wZ)^2+D_wB_w^2-D_uD_wZ^2.
\end{gathered}
\]
The target is exactly
\[
\begin{split}
 F_w={}&4mM_w^2+4CD_wR_w^2+4AD_uD_wY^2+4DJ_w^2\\
       &+4\alpha D_u(Y-wZ)^2+4\beta D_wB_w^2+2K_w.
\end{split}                                                     \tag{P4}
\]
This expression, not a restriction to a contact line, is preserved
throughout the proof. The source and testing angles are never identified.

\paragraph{Why a strict bounded Gram suffices.}
The only lifting lemma used below is the established separated
bounded-Gram lemma: a finite symmetric matrix whose entries are
finite sums of bounded separated factors, and whose real residue
is strictly positive definite, defines a tensor SOS. Its elementary
mechanism is worth recording. For a source infinitesimal $s$ and
bounded testing $t$, choose positive real $a,\varepsilon$ with
$|t|<a$ and $|s|<\varepsilon/a$. Then exactly
\[
 \varepsilon\pm st=\tfrac12
 \{(\varepsilon/a+s)(a\pm t)
            +(\varepsilon/a-s)(a\mp t)\}.
\]
Each summand is a product of nonnegative pure-field factors and
has a separated square root. The same argument interchanges the
fields. A finite real diagonal budget absorbs all infinitesimal
matrix entries into a strictly positive real Gram. This constructs
finite SOS without a mixed Cholesky factor or inverse. Every later
use explicitly supplies the bounded coefficients and the positive
residue; symbolic checks alone do not establish this hypothesis.

\subsection*{The exact complete normal identity used in the closure}

Here is the complete SP6 identity, with the centered $A,\alpha$
excesses restored, so that no cap term is implicit. Let
\[
\begin{gathered}
 q=\sqrt z,\quad t=u^2,\ \xi=w^2,\quad g=h-Q,
 \quad b_\theta=\beta\theta,\quad a_\beta=\beta\theta^2,\\
 N_x=X+uhZ,\quad N_y=(Y-bwZ)/q,\\
 \Gamma=\sigma Z-uwW,\quad R=wZ/q-uqW,
 \quad A_\theta=q(t-x)W,\\
 P_\beta=uZ-(x/\sigma)wW,\quad X_e=X+ewW,\quad
 V=uw(t-x)W.
\end{gathered}                                                   \tag{P5}
\]
The source coefficients are
\[
\begin{split}
 \nu&=(2D-2b_\theta+h)/\sigma,\\
 \lambda&=(8Dh+4a_\beta+2h^2)/\sigma-4bp\sigma z
            =2H+\lambda_0,\\
 \lambda_0&=4D(1+h)/\sigma-2\theta\nu-4bp\sigma z,\\
 \chi&=8(pb+lQ)-4h+4eh(2D+1)-8Dh\\
      &\quad+4z\{4Dbh+4pb\sigma+(1+2m)g^2+2a_\beta\},\\
 B_x&=8\kappa e,\qquad
 A_y=8(p-l)+4z\{(1+2m)g+2Dh+2p\sigma\},\\
 R_x&=-4\{2Db+(1+2m)g-2b_\theta\},\\
 R_\gamma&=-4\{2(C+D)(r-1)+2p+r-1+2pt\},\\
 R_y&=4\{t[2Dh+2p\sigma+(1+2m)g]
       +b(2C+2D+1)(\xi+\sigma z)\\
      &\hspace{22mm}+(2C+1)(b-\sigma)+2D(1-\sigma)+2p\sigma\}.
\end{split}                                                     \tag{P6}
\]
Finally put
\[
\begin{split}
 k_0&=2b\{(2C+1)b+4D+4p\sigma
                 +2(2C+2D+1)\sigma(r-1)\},\\
 k_1&=2(1+2m)g^2+8Dhb+8pb\sigma+4a_\beta,\\
 k_2&=2b^2(2C+2D+1),\qquad K_r=k_0+k_1t+k_2\xi,\\
 H_y&=4D_u(aD_w+\ell)+4mtz+4CzD_w+(4D+2)\xi z+2D_uz,\\
 c_G&=4D_u(pbD_w+lQ),\qquad c_N=8D_u(pD_w-l).
\end{split}                                                     \tag{P7}
\]
Then $F_w=4\beta D_wN_x^2+G$, with the full parent core
\[
\begin{split}
 G={}&2(2D+1)(uX_e)^2+[2(1+\xi)+4m\xi]N_x^2+H_yN_y^2\\
 &+(4D/\sigma^2+4bpzt)\Gamma^2+(c_G+zK_r)R^2
                -4u\nu N_x\Gamma\\
 &-4uqw(2m+2D+1)N_xN_y
      +qwR_\gamma(N_y+bR)\Gamma\\
 &+uqwR_xN_xR+(c_N+zR_y)N_yR
           +qwB_xN_xA_\theta+uA_yN_yA_\theta\\
 &+4\kappa H A_\theta^2
       +\lambda u(\Gamma P_\beta+A_\theta R)+\chi uA_\theta R.
\end{split}                                                     \tag{SP6}
\]
This is an exact polynomial identity after clearing pure-source
denominators in P1. It is the imported JR5/JB1 identity followed
by the explicitly checked SP1--SP6 transformation. P4--P7 specify
every entry for direct hand reproduction; the complete earlier
derivations and their checkers are preserved in the archive.
The two crucial conductors are
\[
\begin{split}
 A_\theta^2+u\Gamma P_\beta+uA_\theta R
  &=\sigma P_\beta^2+(x/\sigma)(t-x)(\xi-\sigma z)W^2,\\
 V^2&=\sigma tA_\theta^2+u\sigma(t-x)A_\theta R-(t-x)V\Gamma.
\end{split}                                                     \tag{P8}
\]
For the bounded-$M_\beta$ case, the generic pre-extraction JR
coefficients of $\Gamma^2,A_\theta^2,\Gamma V$ are, respectively,
\[
\begin{split}
 Z_0&=\{4D+2t(4Dh+2a_\beta+h^2)
                       +2t^2h^2(2D+1)\}/\sigma^2,\\
 A_0&=4\kappa H+2h^2(2D+1)t/\sigma,\\
 S&=2\{h^2(2D+1)(t+x)+4Dh-2bp\sigma^2z
                                  +2a_\beta+h^2\}/\sigma^2.
\end{split}
\]
Its other changed terms are $H_xN_x^2+Z_xN_x\Gamma+uA_xN_xV
+u\sigma B_gA_\theta R$, where
\[
\begin{split}
 H_x&=2(1+t+\xi)+4Dt+4m\xi,\quad A_x=-4h(2D+1)/\sigma,\\
 Z_x&=-4u\{2D-2b_\theta+h+th(2D+1)\}/\sigma,\\
 \sigma B_g&=2h^2(2D+1)(t-x)/\sigma+\lambda+\chi.
\end{split}
\]
These replace the corresponding extracted terms in SP6; all other
normal and radial terms are unchanged.

\paragraph{The needed pole cancellations.}
Directly from P1--P2,
\[
 \nu=2D/\sigma+\{H_0-ez(J-\sigma)\}/P.
\]
To bound $\alpha_\kappa=2H(2\kappa-1)$, set
$N_\kappa=K(1+\sigma r)(1+\sigma z)-MP$ and
\[
 F_x=N_\kappa|_{\sigma=0}/x,\qquad
 F_\sigma=(N_\kappa-N_\kappa|_{\sigma=0})/\sigma.
\]
These quotients are polynomials in bounded source letters, and
$\alpha_\kappa=2(hF_\sigma+eF_x)/(MP)$. This establishes
boundedness through every opening before taking residues.
The divergent $C$ block is treated by CT7--CT7a. Equivalently,
for $J_C=N_y+bR$ it is exactly
\[
 4CzD_wJ_C^2+8C(r-1)J_C(z\sigma R-qw\Gamma).
\]
The optional cap $C_*=C/(1+C)$ therefore changes the entire core
by precisely $4zD_w(C-C_*)J_C^2
-8(C-C_*)(1-r)J_C(z\sigma R-qw\Gamma)$, as used in HG and HP.
No omitted matrix entry is hidden by that cap.

